\chapter{The Fuzzy Ideal Monad, Fuzzy Domains, and Fuzzy Scott Opens}
\label{ch:fuzzy-ideal-monad-domains-scott}

In the previous chapters we developed fuzzy predicates as enriched presheaves and
then localized their predicate frames to obtain fuzzy frames and fuzzy locales.  We
now apply the same predicate construction to \(X^{\op}\).  Thus the upper fuzzy
predicate object of a small \(\mathbb V\)-category \(X\) is
\[
    P_{\mathbb V}(X^{\op})
    =
    [X,\underline{\mathbb V}],
\]
and, in the posetal specialization determined by
\[
    \Lambda=(L,\leq,\meet,\join,\tensor,\multimap,\bot_L,\top_L),
\]
it is
\[
    P_\Lambda(X^{\op})
    =
    [X,\Lambda].
\]
Its elements are upper fuzzy predicates on \(X\).

The central construction of this chapter is the character adjunction
\[
    \mathcal P_{\mathbb V}^{\uparrow}
    \dashv
    \operatorname{Char}_{\mathbb V}.
\]
Here \(\mathcal P_{\mathbb V}^{\uparrow}X\) is the upper fuzzy Alexandrov locale
whose frame of opens is \([X,\underline{\mathbb V}]\), while
\(\operatorname{Char}_{\mathbb V}\) sends a fuzzy locale to its
\(\mathbb V\)-category of fuzzy-frame morphisms into the truth-value fuzzy frame
\(\underline{\mathbb V}\).  Thus characters do not target the fuzzy Sierpiński
frame.  The fuzzy Sierpiński frame classifies fuzzy opens; the truth-value frame
receives fuzzy characters.

The induced monad on fuzzy preorders is the fuzzy ideal monad:
\[
    T_{\mathbb V}X
    \simeq
    \operatorname{Idl}_{\mathbb V}(X),
\]
where \(\operatorname{Idl}_{\mathbb V}(X)\) is the \(\mathbb V\)-category of
finite-flat weights
\[
    X^{\op}\to\underline{\mathbb V}.
\]
In the posetal specialization,
\[
    T_\Lambda X
    \cong
    \operatorname{Idl}_\Lambda(X),
\]
whose elements are fuzzy ideals of \(X\).

A weak fuzzy domain is a fuzzy preorder equipped with a chosen operation
\[
    \alpha:\operatorname{Idl}_\Lambda(X)\to X
\]
assigning candidate joins to fuzzy ideals.  A fuzzy domain is a strict algebra for
the fuzzy ideal monad.  In the crisp case \(\Lambda=2\), this recovers the ordinary
ideal-completion monad, and strict algebras on separated preorders are precisely
dcpos.

The same monad also contains the algebraic and continuous layers of domain theory.
Free fuzzy domains
\[
    \operatorname{Idl}_\Lambda(B)
\]
are the presented fuzzy algebraic domains.  Fuzzy continuous domains are fuzzy
domains for which the ideal-join map
\[
    \alpha:\operatorname{Idl}_\Lambda(X)\to X
\]
admits an enriched left adjoint
\[
    \beta:X\to\operatorname{Idl}_\Lambda(X).
\]
In the crisp case, \(\beta(x)\) is the ideal of elements way-below \(x\).

Finally, fuzzy Scott opens are upper fuzzy predicates preserving fuzzy ideal joins.
For a weak fuzzy domain
\[
    \alpha:\operatorname{Idl}_\Lambda(X)\to X,
\]
a fuzzy Scott open is an upper fuzzy predicate
\[
    U:X\to L
\]
satisfying
\[
    U(\alpha(I))
    =
    \bigjoin_{x\in X} I(x)\tensor U(x)
\]
for every fuzzy ideal \(I\).

\begin{remark}[Universe and doctrine conventions]
\label{rem:ideal-monad-universe-and-doctrine-convention}
We use the same universe convention as in the preceding chapters.  If
\(\mathcal A\) is not small, then
\[
    [\mathcal A,\underline{\mathbb V}]
    \qquad\text{and}\qquad
    \operatorname{Char}_{\mathbb V}(\mathcal A)
\]
are understood in the appropriate larger universe.  In particular, for small
\(X\), the category
\[
    T_{\mathbb V}X
    =
    \operatorname{Char}_{\mathbb V}(P_{\mathbb V}(X^{\op}))
\]
is taken in the universe where the corresponding category of finite-flat weights
exists.

At the categorical-cosmos level, \(T_{\mathbb V}\) is therefore a
universe-relative monad on the chosen \(\mathbb V\)-CAT of locally small
\(\mathbb V\)-categories for which the finite-flat-weight category exists in the
ambient universe.  In the small posetal specialization it restricts to an honest
monad
\[
    T_\Lambda:\Pre_\Lambda\to\Pre_\Lambda.
\]

Throughout the chapter, ``finite-flat'' means flat with respect to the chosen
finite-limit doctrine used in the definition of categorical fuzzy frames.  In the
posetal scalar-frame specialization, this doctrine consists of the terminal element
and binary meets.  Thus finite-flatness specializes to preservation of top and
binary meets.
\end{remark}

\begin{remark}[Strictification of the ideal monad]
\label{rem:strictification-ideal-monad}
The enriched character construction gives
\[
    T_{\mathbb V}X
    =
    \operatorname{Char}_{\mathbb V}(P_{\mathbb V}(X^{\op}))
    \simeq
    \operatorname{Idl}_{\mathbb V}(X).
\]
After the finite-flat-weight representation is established, we freely replace this
monad by the transported equivalent monad whose value on \(X\) is
\(\operatorname{Idl}_{\mathbb V}(X)\), with unit and multiplication given by the
Yoneda and flattening formulas below.  In the posetal specialization this
identification is an isomorphism of fuzzy preorders.
\end{remark}

\section{Upper fuzzy predicates and upper fuzzy Alexandrov locales}
\label{sec:upper-fuzzy-predicates}

\subsection{The enriched construction}
\label{subsec:upper-predicates-enriched}

\begin{definition}[Upper fuzzy predicates]
\label{def:upper-fuzzy-predicates-enriched}
Let \(X\) be a small \(\mathbb V\)-category.  The \textbf{upper fuzzy predicate
category} of \(X\) is
\[
    P_{\mathbb V}(X^{\op})
    =
    [X,\underline{\mathbb V}].
\]
Its objects are \(\mathbb V\)-functors
\[
    U:X\to\underline{\mathbb V}.
\]
\end{definition}

\begin{definition}[Upper fuzzy Alexandrov locale]
\label{def:upper-fuzzy-alexandrov-locale}
The \textbf{upper fuzzy Alexandrov locale} associated to \(X\) is the fuzzy locale
\[
    \mathcal P_{\mathbb V}^{\uparrow}X
\]
whose fuzzy frame of opens is
\[
    \operatorname{Opens}_{\mathbb V}(\mathcal P_{\mathbb V}^{\uparrow}X)
    =
    P_{\mathbb V}(X^{\op})
    =
    [X,\underline{\mathbb V}].
\]
Since \(P_{\mathbb V}(X^{\op})\) is itself a fuzzy presheaf category, it is a
categorical fuzzy frame via the identity left-exact localization
\[
    1\dashv 1:
    P_{\mathbb V}(X^{\op})
    \rightleftarrows
    P_{\mathbb V}(X^{\op}).
\]
Thus it defines a fuzzy locale by duality.
\end{definition}

\begin{proposition}[Functoriality of upper fuzzy predicates]
\label{prop:upper-fuzzy-predicates-functorial}
Let
\[
    f:X\to Y
\]
be a \(\mathbb V\)-functor.  Precomposition defines a fuzzy-frame morphism
\[
    f^\ast:
    P_{\mathbb V}(Y^{\op})
    =
    [Y,\underline{\mathbb V}]
    \longrightarrow
    [X,\underline{\mathbb V}]
    =
    P_{\mathbb V}(X^{\op}),
    \qquad
    f^\ast(V)=Vf.
\]
Consequently there is a fuzzy-locale map
\[
    \mathcal P_{\mathbb V}^{\uparrow}X
    \longrightarrow
    \mathcal P_{\mathbb V}^{\uparrow}Y,
\]
and hence a covariant functor
\[
    \mathcal P_{\mathbb V}^{\uparrow}:
    \mathbb V\text{-}\mathbf{Cat}
    \longrightarrow
    \mathbf{FLoc}_{\mathbb V}.
\]
\end{proposition}

\begin{proof}
All fuzzy-frame operations in the predicate frame are computed pointwise.
Precomposition preserves them.  Passing to frame opposites gives the displayed
locale map.
\end{proof}

\subsection{The posetal specialization}
\label{subsec:upper-predicates-posetal}

\begin{definition}[Upper fuzzy predicate over \(\Lambda\)]
\label{def:upper-fuzzy-predicate-posetal}
Let \(X\) be a fuzzy preorder over \(\Lambda\).  An \textbf{upper fuzzy predicate}
on \(X\) is a map
\[
    U:X\to L
\]
satisfying
\[
    X(x,x')\tensor U(x)\leq U(x')
\]
for all \(x,x'\in X\).
\end{definition}

\begin{proposition}[The upper Alexandrov fuzzy frame]
\label{prop:upper-alexandrov-fuzzy-frame}
The upper fuzzy predicates on \(X\) form the fuzzy frame
\[
    P_\Lambda(X^{\op})
    =
    [X,\Lambda].
\]
Its hom-values are
\[
    P_\Lambda(X^{\op})(U,V)
    =
    \bigmeet_{x\in X}
    \bigl(U(x)\multimap V(x)\bigr),
\]
and its induced crisp order is pointwise.  Joins, meets, and scalar multiplication
are also computed pointwise:
\[
    \left(\bigjoin_i U_i\right)(x)=\bigjoin_i U_i(x),
    \qquad
    \left(\bigmeet_i U_i\right)(x)=\bigmeet_i U_i(x),
    \qquad
    (r\tensor U)(x)=r\tensor U(x).
\]
\end{proposition}

\begin{proof}
This is the predicate-frame construction from the previous chapters applied to
\(X^{\op}\).
\end{proof}

\begin{example}[The crisp case]
\label{ex:upper-predicates-crisp}
For \(\Lambda=2\), a fuzzy preorder is an ordinary preorder and
\[
    P_2(X^{\op})
    \cong
    \operatorname{Up}(X).
\]
Indeed, the upper-predicate law becomes
\[
    x\leq x'
    \quad\text{and}\quad
    U(x)
    \quad\Longrightarrow\quad
    U(x').
\]
Later, Theorem~\ref{thm:scott-opens-free-fuzzy-domain} identifies this frame with
the Scott frame of the free dcpo \(\operatorname{Idl}(X)\).
\end{example}

\begin{example}[Standard truth-value objects]
\label{ex:standard-upper-fuzzy-predicates}
For the G\"odel tensor \(r\tensor s=r\meet s\), an upper fuzzy predicate satisfies
\[
    X(x,x')\meet U(x)\leq U(x').
\]
For the {\L}ukasiewicz tensor \(r\tensor s=\max(0,r+s-1)\), it satisfies
\[
    \max(0,X(x,x')+U(x)-1)\leq U(x').
\]
For the product tensor \(r\tensor s=rs\), it satisfies
\[
    X(x,x')\,U(x)\leq U(x').
\]
For the Lawvere quantale \([0,\infty]\), with the Lawvere order
\(a\leq_\Lambda b\) iff \(a\geq b\) in the usual order, tensor \(+\), top \(0\),
and bottom \(\infty\), the upper-predicate law becomes
\[
    U(x')\leq d(x,x')+U(x)
\]
in the usual numerical order.
\end{example}

\section{Fuzzy characters}
\label{sec:fuzzy-characters}

\subsection{The enriched construction}
\label{subsec:fuzzy-characters-enriched}

The truth-value object \(\underline{\mathbb V}\) is itself a categorical fuzzy
frame.  Indeed, if \(\mathbf 1\) denotes the terminal small \(\mathbb V\)-category,
then
\[
    P_{\mathbb V}(\mathbf 1)
    =
    [\mathbf 1^{\op},\underline{\mathbb V}]
    \simeq
    \underline{\mathbb V}.
\]
Thus \(\underline{\mathbb V}\) is the presheaf fuzzy frame on \(\mathbf 1\).

\begin{definition}[Fuzzy character]
\label{def:fuzzy-character-enriched}
Let \(\mathcal A\) be a categorical fuzzy frame.  Its \textbf{fuzzy character
category} is the full \(\mathbb V\)-subcategory
\[
    \operatorname{Char}_{\mathbb V}(\mathcal A)
    \subseteq
    [\mathcal A,\underline{\mathbb V}]
\]
spanned by those \(\mathbb V\)-functors
\[
    \chi:\mathcal A\to\underline{\mathbb V}
\]
that are fuzzy-frame morphisms, that is, that preserve colimits, the chosen finite
limits, and scalars.
\end{definition}

\begin{definition}[Characters of a fuzzy locale]
\label{def:fuzzy-characters-locale}
If \(\mathcal X\) is a fuzzy locale, define
\[
    \operatorname{Char}_{\mathbb V}(\mathcal X)
    \coloneqq
    \operatorname{Char}_{\mathbb V}
    (\operatorname{Opens}_{\mathbb V}(\mathcal X)).
\]
\end{definition}

\begin{remark}[Presentation-wise characters]
\label{rem:character-universe-convention}
If a fuzzy frame is presented as a reflective localization
\[
    a\dashv i:
    P_{\mathbb V}(X)\rightleftarrows\mathcal A,
\]
then characters of \(\mathcal A\) may be described as characters of
\(P_{\mathbb V}(X)\) satisfying the descent condition imposed by the localization.
For upper predicate frames, this descent-free case becomes the finite-flat-weight
description of Section~\ref{sec:fuzzy-ideals-flat-weights}.
\end{remark}

\subsection{The posetal specialization}
\label{subsec:fuzzy-characters-posetal}

\begin{definition}[Fuzzy character over \(\Lambda\)]
\label{def:fuzzy-character-posetal}
Let \(A\in\FFrm_\Lambda\).  Its \textbf{fuzzy character preorder}
\[
    \operatorname{Char}_\Lambda(A)
\]
is the full fuzzy subpreorder of \([A,\Lambda]\) spanned by the scalar-frame
homomorphisms
\[
    \chi:A\to\Lambda.
\]
Equivalently, its underlying set is
\[
    |\operatorname{Char}_\Lambda(A)|
    =
    \FFrm_\Lambda(A,\Lambda).
\]
Thus an element of \(\operatorname{Char}_\Lambda(A)\) is a map preserving arbitrary
joins, finite meets, and scalar multiplication:
\[
    \chi(r\tensor_A a)=r\tensor\chi(a).
\]
\end{definition}

\begin{proposition}[The fuzzy preorder of characters]
\label{prop:fuzzy-character-preorder}
The fuzzy preorder on \(\operatorname{Char}_\Lambda(A)\) is inherited from the
pointwise functor preorder:
\[
    \operatorname{Char}_\Lambda(A)(\chi,\psi)
    =
    \bigmeet_{a\in A}
    \bigl(\chi(a)\multimap\psi(a)\bigr).
\]
Its induced crisp order is pointwise.
\end{proposition}

\begin{proof}
The character preorder is a full fuzzy subpreorder of \([A,\Lambda]\).
\end{proof}

\begin{definition}[Characters of a fuzzy locale over \(\Lambda\)]
\label{def:fuzzy-character-locale-posetal}
For a fuzzy locale \(\mathcal X\), define
\[
    \operatorname{Char}_\Lambda(\mathcal X)
    \coloneqq
    \operatorname{Char}_\Lambda(\Opens_\Lambda(\mathcal X)).
\]
\end{definition}

\begin{proposition}[Characters of the truth-value frame]
\label{prop:characters-truth-value-frame}
There is exactly one scalar-frame character
\[
    \Lambda\to\Lambda,
\]
namely the identity.  Hence
\[
    \operatorname{Char}_\Lambda(\Lambda)
\]
is the terminal fuzzy preorder.
\end{proposition}

\begin{proof}
Let \(\chi:\Lambda\to\Lambda\) be a scalar-frame homomorphism.  Since \(\chi\)
preserves top,
\[
    \chi(\top_L)=\top_L.
\]
For every \(r\in L\),
\[
    \chi(r)
    =
    \chi(r\tensor\top_L)
    =
    r\tensor\chi(\top_L)
    =
    r\tensor\top_L
    =
    r.
\]
Thus \(\chi=1_\Lambda\).
\end{proof}

\begin{example}[The crisp case]
\label{ex:fuzzy-characters-crisp}
For \(\Lambda=2\), a character of a frame \(A\) is a frame homomorphism
\[
    A\to 2,
\]
equivalently a completely prime filter of \(A\).  Thus, in the crisp localic case,
characters are points.
\end{example}

\section{The upper fuzzy character adjunction}
\label{sec:upper-fuzzy-character-adjunction}

\subsection{The enriched adjunction}
\label{subsec:upper-character-adjunction-enriched}

\begin{theorem}[Upper fuzzy character adjunction]
\label{thm:upper-fuzzy-character-adjunction-enriched}
Let \(X\) be a small \(\mathbb V\)-category and let \(\mathcal X\) be a fuzzy
locale.  There is a natural equivalence
\[
    \mathbf{FLoc}_{\mathbb V}
    (\mathcal P_{\mathbb V}^{\uparrow}X,\mathcal X)
    \simeq
    \mathbb V\text{-}\mathbf{Cat}
    \bigl(
        X,
        \operatorname{Char}_{\mathbb V}(\mathcal X)
    \bigr).
\]
Equivalently, writing
\[
    A=\operatorname{Opens}_{\mathbb V}(\mathcal X),
\]
there is a natural equivalence
\[
    \mathbf{FFrm}_{\mathbb V}
    \bigl(
        A,
        P_{\mathbb V}(X^{\op})
    \bigr)
    \simeq
    \mathbb V\text{-}\mathbf{Cat}
    \bigl(
        X,
        \operatorname{Char}_{\mathbb V}(A)
    \bigr).
\]
\end{theorem}

\begin{proof}
A fuzzy-locale map
\[
    \mathcal P_{\mathbb V}^{\uparrow}X\to\mathcal X
\]
is frame-dually a fuzzy-frame morphism
\[
    h:A\to [X,\underline{\mathbb V}].
\]
By the enriched tensor--hom adjunction and the symmetry of \(\tensor\), this is
equivalently a \(\mathbb V\)-functor
\[
    \widehat h:X\to [A,\underline{\mathbb V}],
    \qquad
    \widehat h(x)(a)=h(a)(x).
\]
Explicitly, this is the composite transposition
\[
    \mathbb V\text{-}\mathbf{CAT}(A,[X,\underline{\mathbb V}])
    \cong
    \mathbb V\text{-}\mathbf{CAT}(A\tensor X,\underline{\mathbb V})
    \cong
    \mathbb V\text{-}\mathbf{CAT}(X,[A,\underline{\mathbb V}]).
\]

Since evaluation at \(x\) preserves the pointwise fuzzy-frame structure on
\([X,\underline{\mathbb V}]\), each
\[
    \widehat h(x):A\to\underline{\mathbb V}
\]
is a fuzzy character of \(A\).  Hence \(\widehat h\) factors through
\[
    \operatorname{Char}_{\mathbb V}(A)
    \subseteq
    [A,\underline{\mathbb V}].
\]

Conversely, a \(\mathbb V\)-functor
\[
    k:X\to\operatorname{Char}_{\mathbb V}(A)
\]
defines
\[
    \widetilde k:A\to[X,\underline{\mathbb V}],
    \qquad
    \widetilde k(a)(x)=k(x)(a).
\]
For each \(a\), this is an upper predicate on \(X\); for each \(x\), the map
\(a\mapsto k(x)(a)\) preserves the fuzzy-frame structure.  Therefore
\(\widetilde k\) is a fuzzy-frame morphism.  The two constructions are inverse by
the enriched tensor--hom adjunction, and naturality is inherited from the
naturality of that adjunction.
\end{proof}

\subsection{The posetal specialization}
\label{subsec:upper-character-adjunction-posetal}

\begin{theorem}[Posetal upper fuzzy character adjunction]
\label{thm:upper-fuzzy-character-adjunction-posetal}
Let \(X\) be a fuzzy preorder and let \(A\in\FFrm_\Lambda\).  There is a natural
isomorphism
\[
    \FFrm_\Lambda(A,P_\Lambda(X^{\op}))
    \cong
    \Pre_\Lambda(X,\operatorname{Char}_\Lambda(A)).
\]
\end{theorem}

\begin{proof}
For a fuzzy-frame homomorphism
\[
    h:A\to P_\Lambda(X^{\op}),
\]
define
\[
    \widehat h(x)(a)=h(a)(x).
\]
Each \(\widehat h(x)\) is a character because evaluation at \(x\) preserves the
pointwise scalar-frame structure.  Moreover, for all \(x,x'\in X\) and \(a\in A\),
the upper-predicate law for \(h(a)\) gives
\[
    X(x,x')\tensor h(a)(x)\leq h(a)(x'),
\]
or equivalently
\[
    X(x,x')\leq h(a)(x)\multimap h(a)(x').
\]
Taking the meet over \(a\) shows that
\[
    X(x,x')
    \leq
    \operatorname{Char}_\Lambda(A)(\widehat h(x),\widehat h(x')).
\]
Thus \(\widehat h\) is a \(\Lambda\)-functor.

Conversely, a \(\Lambda\)-functor
\[
    k:X\to\operatorname{Char}_\Lambda(A)
\]
defines
\[
    \widetilde k(a)(x)=k(x)(a).
\]
Functoriality of \(k\) gives the upper-predicate law in \(x\), while the character
condition on each \(k(x)\) gives preservation of joins, finite meets, and scalars
in \(a\).  Hence \(\widetilde k\) is a fuzzy-frame homomorphism.  The two
assignments are inverse.
\end{proof}

\begin{example}[The crisp case]
\label{ex:upper-character-adjunction-crisp}
For \(\Lambda=2\), the adjunction becomes
\[
    \Frm(A,\operatorname{Up}(X))
    \cong
    \Pre(X,\operatorname{Char}_2(A)).
\]
Thus frame maps into an upper Alexandrov frame are the same as monotone
point-valued families indexed by \(X\).
\end{example}

\begin{example}[Discrete fuzzy preorders]
\label{ex:upper-character-adjunction-discrete}
If \(X\) is discrete, then
\[
    P_\Lambda(X^{\op})\cong L^X.
\]
A fuzzy-frame homomorphism
\[
    A\to L^X
\]
is precisely an \(X\)-indexed family of characters \(A\to\Lambda\).
\end{example}

\section{Fuzzy ideals as finite-flat weights}
\label{sec:fuzzy-ideals-flat-weights}

\subsection{The enriched construction}
\label{subsec:fuzzy-ideals-enriched}

\begin{definition}[Finite-flat weight]
\label{def:finite-flat-weight}
Let \(X\) be a small \(\mathbb V\)-category.  A \textbf{finite-flat weight} on
\(X\) is a weight
\[
    I:X^{\op}\to\underline{\mathbb V}
\]
such that the weighted-colimit functional
\[
    \chi_I:[X,\underline{\mathbb V}]\to\underline{\mathbb V},
    \qquad
    \chi_I(U)=\int^{x\in X} I(x)\tensor U(x),
\]
preserves the chosen finite limits of the fuzzy-frame doctrine.  Equivalently,
\(I\)-weighted colimits in \(\underline{\mathbb V}\) commute with those finite
limits.
\end{definition}

\begin{definition}[Fuzzy ideal]
\label{def:fuzzy-ideal-enriched}
A \textbf{fuzzy ideal} of a small \(\mathbb V\)-category \(X\) is a finite-flat
weight
\[
    I:X^{\op}\to\underline{\mathbb V}.
\]
We write
\[
    \operatorname{Idl}_{\mathbb V}(X)
\]
for the full \(\mathbb V\)-subcategory of
\[
    [X^{\op},\underline{\mathbb V}]
\]
spanned by the finite-flat weights.
\end{definition}

\begin{proposition}[Flat-weight representation of cocontinuous functionals]
\label{prop:flat-weight-representation-of-characters}
Let \(X\) be small.  Cocontinuous \(\mathbb V\)-functors
\[
    [X,\underline{\mathbb V}]\to\underline{\mathbb V}
\]
are, up to canonical isomorphism, precisely the weighted-colimit functionals
\[
    U\longmapsto \int^{x\in X} I(x)\tensor U(x)
\]
for weights
\[
    I:X^{\op}\to\underline{\mathbb V}.
\]
Such a functional preserves the chosen finite limits precisely when \(I\) is
finite-flat.
\end{proposition}

\begin{proof}
Let
\[
    \chi:[X,\underline{\mathbb V}]\to\underline{\mathbb V}
\]
be cocontinuous.  Define
\[
    I_\chi(x)=\chi(X(x,-)).
\]
The assignment
\[
    x\longmapsto X(x,-)
\]
is the Yoneda embedding of \(X^{\op}\) into \([X,\underline{\mathbb V}]\).  Since
\(\chi\) is a \(\mathbb V\)-functor, \(I_\chi\) is a weight
\[
    I_\chi:X^{\op}\to\underline{\mathbb V}.
\]

For every
\[
    U:X\to\underline{\mathbb V},
\]
the enriched co-Yoneda formula in the copresheaf category
\([X,\underline{\mathbb V}]\) gives
\[
    U
    \cong
    \int^{x\in X} U(x)\tensor X(x,-).
\]
Applying the cocontinuous functor \(\chi\), we obtain
\[
\begin{aligned}
    \chi(U)
    &\cong
    \chi\left(\int^{x\in X} U(x)\tensor X(x,-)\right)\\
    &\cong
    \int^{x\in X} U(x)\tensor\chi(X(x,-))\\
    &=
    \int^{x\in X} U(x)\tensor I_\chi(x)\\
    &\cong
    \int^{x\in X} I_\chi(x)\tensor U(x).
\end{aligned}
\]
Thus \(\chi\) is represented by \(I_\chi\).

Conversely, every weight \(I:X^{\op}\to\underline{\mathbb V}\) determines the
weighted-colimit functional
\[
    \chi_I(U)=\int^{x\in X}I(x)\tensor U(x).
\]
This functional is cocontinuous because it is built from tensors and small
colimits in \(\underline{\mathbb V}\), and tensors preserve colimits in each
variable.  The two constructions are inverse by the same co-Yoneda calculation,
since
\[
    \chi_I(X(x,-))
    \cong
    \int^{y\in X}I(y)\tensor X(x,y)
    \cong
    I(x).
\]

Finally, preservation of the chosen finite limits by \(\chi_I\) is exactly the
definition of finite-flatness of the weight \(I\).
\end{proof}

\begin{theorem}[Characters of upper predicate frames]
\label{thm:characters-upper-predicate-frames-flat}
For every small \(\mathbb V\)-category \(X\), there is a canonical equivalence
\[
    \operatorname{Char}_{\mathbb V}(P_{\mathbb V}(X^{\op}))
    \simeq
    \operatorname{Idl}_{\mathbb V}(X).
\]
Under this equivalence, a fuzzy ideal
\[
    I:X^{\op}\to\underline{\mathbb V}
\]
corresponds to the character
\[
    \chi_I(U)
    =
    \int^{x\in X}I(x)\tensor U(x).
\]
\end{theorem}

\begin{proof}
Since
\[
    P_{\mathbb V}(X^{\op})=[X,\underline{\mathbb V}],
\]
a character is a cocontinuous, finite-limit-preserving, scalar-preserving
\(\mathbb V\)-functor
\[
    [X,\underline{\mathbb V}]\to\underline{\mathbb V}.
\]
By Proposition~\ref{prop:flat-weight-representation-of-characters}, the
cocontinuous finite-limit-preserving functors are exactly the finite-flat
weighted-colimit functionals
\[
    U\longmapsto\int^{x\in X}I(x)\tensor U(x).
\]
Scalar preservation is automatic for these functionals because tensoring
distributes over the coends defining the weighted colimits:
\[
\begin{aligned}
    \chi_I(r\tensor U)
    &=
    \int^{x\in X} I(x)\tensor r\tensor U(x)\\
    &\cong
    r\tensor
    \int^{x\in X}I(x)\tensor U(x)\\
    &=
    r\tensor\chi_I(U).
\end{aligned}
\]
Thus characters of the upper predicate frame are precisely finite-flat weights on
\(X\).
\end{proof}

\subsection{The posetal specialization}
\label{subsec:fuzzy-ideals-posetal}

\begin{definition}[Fuzzy ideal over \(\Lambda\)]
\label{def:fuzzy-ideal-posetal}
Let \(X\) be a fuzzy preorder over \(\Lambda\).  A \textbf{fuzzy ideal} of \(X\) is
a lower fuzzy predicate
\[
    I:X^{\op}\to\Lambda
\]
such that
\[
    \chi_I(U)
    =
    \bigjoin_{x\in X} I(x)\tensor U(x)
\]
preserves top and binary meets as a map
\[
    P_\Lambda(X^{\op})\to\Lambda.
\]
Equivalently, \(I\) satisfies the lower-predicate law
\[
    X(x',x)\tensor I(x)\leq I(x')
\]
and the finite-flatness equations
\[
    \bigjoin_{x\in X}I(x)=\top_L
\]
and
\[
    \bigjoin_{x\in X}I(x)\tensor\bigl(U(x)\meet V(x)\bigr)
    =
    \left(\bigjoin_{x\in X}I(x)\tensor U(x)\right)
    \meet
    \left(\bigjoin_{x\in X}I(x)\tensor V(x)\right)
\]
for all upper fuzzy predicates \(U,V\).
\end{definition}

\begin{definition}[The fuzzy preorder of fuzzy ideals]
\label{def:fuzzy-ideal-preorder-posetal}
For a fuzzy preorder \(X\), define
\[
    \operatorname{Idl}_\Lambda(X)
\]
to be the full fuzzy subpreorder of \(P_\Lambda(X)\) spanned by the fuzzy ideals.
Thus, for fuzzy ideals \(I,J\),
\[
    \operatorname{Idl}_\Lambda(X)(I,J)
    =
    \bigmeet_{x\in X}
    \bigl(I(x)\multimap J(x)\bigr).
\]
\end{definition}

\begin{theorem}[Characters of upper fuzzy Alexandrov frames]
\label{thm:characters-upper-alexandrov-fuzzy-ideals}
For every fuzzy preorder \(X\), there is an isomorphism of fuzzy preorders
\[
    \operatorname{Char}_\Lambda(P_\Lambda(X^{\op}))
    \cong
    \operatorname{Idl}_\Lambda(X).
\]
Under this isomorphism, a fuzzy ideal \(I\) corresponds to the character
\[
    \chi_I(U)
    =
    \bigjoin_{x\in X}I(x)\tensor U(x).
\]
\end{theorem}

\begin{proof}
The coend formula in
Theorem~\ref{thm:characters-upper-predicate-frames-flat} specializes to the
displayed join formula, and finite-flatness specializes to preservation of top and
binary meets.

It remains to identify hom-values.  If
\[
    r\leq \bigmeet_x (I(x)\multimap J(x)),
\]
then \(r\tensor I(x)\leq J(x)\) for every \(x\), and therefore
\[
    r\tensor \chi_I(U)
    =
    r\tensor \bigjoin_x I(x)\tensor U(x)
    \leq
    \bigjoin_x J(x)\tensor U(x)
    =
    \chi_J(U)
\]
for every upper predicate \(U\).  Hence
\[
    r\leq
    \bigmeet_U(\chi_I(U)\multimap\chi_J(U)).
\]

Conversely, evaluate at the upper representable \(X(x,-)\).  The enriched Yoneda
calculation gives
\[
    \chi_I(X(x,-))
    =
    \bigjoin_y I(y)\tensor X(x,y)
    =
    I(x),
\]
and similarly for \(J\).  Hence the character hom is bounded above by
\[
    \bigmeet_x(I(x)\multimap J(x)).
\]
Thus the hom-values agree.
\end{proof}

\begin{example}[The crisp case]
\label{ex:fuzzy-ideals-crisp}
For \(\Lambda=2\), a fuzzy ideal is an ordinary ideal
\[
    I\subseteq X.
\]
Thus \(I\) is a lower set, is inhabited, and is directed.  Therefore
\[
    \operatorname{Char}_2(\operatorname{Up}(X))
    \cong
    \operatorname{Idl}(X).
\]
\end{example}

\begin{example}[Standard truth-value objects]
\label{ex:standard-fuzzy-ideals}
For G\"odel truth values,
\[
    \chi_I(U)=\bigjoin_x \min(I(x),U(x)).
\]
For {\L}ukasiewicz truth values,
\[
    \chi_I(U)
    =
    \bigjoin_x \max(0,I(x)+U(x)-1).
\]
For product truth values,
\[
    \chi_I(U)=\bigjoin_x I(x)U(x).
\]
For the Lawvere quantale, joins in \(\Lambda\) are usual infima and tensor is
addition, so
\[
    \chi_I(U)
    =
    \inf_x\bigl(I(x)+U(x)\bigr).
\]
\end{example}

\section{The fuzzy ideal monad}
\label{sec:fuzzy-ideal-monad}

\subsection{The enriched monad}
\label{subsec:fuzzy-ideal-monad-enriched}

\begin{definition}[Fuzzy ideal monad]
\label{def:fuzzy-ideal-monad-enriched}
The \textbf{fuzzy ideal monad} is the monad
\[
    T_{\mathbb V}
    =
    \operatorname{Char}_{\mathbb V}\mathcal P_{\mathbb V}^{\uparrow}
\]
induced by the adjunction
\[
    \mathcal P_{\mathbb V}^{\uparrow}
    \dashv
    \operatorname{Char}_{\mathbb V}.
\]
Thus
\[
    T_{\mathbb V}X
    =
    \operatorname{Char}_{\mathbb V}(P_{\mathbb V}(X^{\op}))
    \simeq
    \operatorname{Idl}_{\mathbb V}(X).
\]
After the strictification convention of
Remark~\ref{rem:strictification-ideal-monad}, we write simply
\[
    T_{\mathbb V}X=\operatorname{Idl}_{\mathbb V}(X).
\]
\end{definition}

\begin{proposition}[Unit and multiplication]
\label{prop:fuzzy-ideal-monad-unit-multiplication-enriched}
Under the finite-flat-weight description, the unit
\[
    \eta_X:X\to\operatorname{Idl}_{\mathbb V}(X)
\]
is the Yoneda embedding into finite-flat weights:
\[
    \eta_X(x)=X(-,x).
\]
The multiplication
\[
    \mu_X:
    \operatorname{Idl}_{\mathbb V}(\operatorname{Idl}_{\mathbb V}(X))
    \to
    \operatorname{Idl}_{\mathbb V}(X)
\]
is flattening:
\[
    \mu_X(\mathcal I)(x)
    =
    \int^{I\in\operatorname{Idl}_{\mathbb V}(X)}
    \mathcal I(I)\tensor I(x).
\]
\end{proposition}

\begin{proof}
The unit of the character adjunction sends \(x\) to evaluation at \(x\):
\[
    U\longmapsto U(x).
\]
The corresponding weight evaluates covariant representables by
\[
    X(y,-)(x)=X(y,x),
\]
so \(\eta_X(x)=X(-,x)\).  The multiplication is the monad multiplication induced by
the adjunction, transported through the finite-flat-weight representation; it is
the weighted colimit of an ideal of ideals.
\end{proof}

\subsection{The posetal specialization}
\label{subsec:fuzzy-ideal-monad-posetal}

\begin{definition}[The posetal fuzzy ideal monad]
\label{def:fuzzy-ideal-monad-posetal}
In the posetal specialization,
\[
    T_\Lambda(X)=\operatorname{Idl}_\Lambda(X).
\]
The unit is
\[
    \eta_X:X\to\operatorname{Idl}_\Lambda(X),
    \qquad
    \eta_X(x)(y)=X(y,x),
\]
and the multiplication is
\[
    \mu_X:
    \operatorname{Idl}_\Lambda(\operatorname{Idl}_\Lambda(X))
    \to
    \operatorname{Idl}_\Lambda(X),
\]
where
\[
    \mu_X(\mathcal I)(x)
    =
    \bigjoin_{I\in\operatorname{Idl}_\Lambda(X)}
    \mathcal I(I)\tensor I(x).
\]
\end{definition}

\begin{proposition}[Functorial action of the fuzzy ideal monad]
\label{prop:fuzzy-ideal-monad-functorial-action}
Let
\[
    f:X\to Y
\]
be a \(\Lambda\)-functor.  Then
\[
    T_\Lambda(f):
    \operatorname{Idl}_\Lambda(X)\to\operatorname{Idl}_\Lambda(Y)
\]
is given by
\[
    T_\Lambda(f)(I)(y)
    =
    \bigjoin_{x\in X}I(x)\tensor Y(y,f(x)).
\]
This is the fuzzy lower closure of the direct image of \(I\) along \(f\).
\end{proposition}

\begin{proof}
This is the posetal specialization of enriched left Kan extension of a weight
along \(f\).  Since it is the functorial action of the monad induced by the
character adjunction, it preserves fuzzy ideals and is functorial in \(f\).
\end{proof}

\begin{example}[The crisp ideal monad]
\label{ex:crisp-ideal-monad}
For \(\Lambda=2\),
\[
    T_2(X)=\operatorname{Idl}(X).
\]
The unit sends
\[
    x\longmapsto\downarrow x,
\]
multiplication sends an ideal of ideals to its union,
\[
    \mu_X(\mathcal I)=\bigcup_{I\in\mathcal I}I,
\]
and for a monotone map \(f:X\to Y\),
\[
    T_2(f)(I)=\downarrow f[I].
\]
Thus \(T_2\) is the ordinary ideal-completion monad.
\end{example}

\begin{example}[Standard direct-image formulas]
\label{ex:standard-ideal-direct-image}
For G\"odel truth values,
\[
    T_\Lambda(f)(I)(y)
    =
    \bigjoin_x \min(I(x),Y(y,f(x))).
\]
For {\L}ukasiewicz truth values,
\[
    T_\Lambda(f)(I)(y)
    =
    \bigjoin_x \max(0,I(x)+Y(y,f(x))-1).
\]
For product truth values,
\[
    T_\Lambda(f)(I)(y)
    =
    \bigjoin_x I(x)Y(y,f(x)).
\]
For the Lawvere quantale,
\[
    T_\Lambda(f)(I)(y)
    =
    \inf_x\bigl(I(x)+d_Y(y,f(x))\bigr).
\]
\end{example}

\section{Weak fuzzy domains and fuzzy domains}
\label{sec:weak-fuzzy-domains-fuzzy-domains}

\subsection{The enriched notion}
\label{subsec:weak-fuzzy-domains-enriched}

\begin{definition}[Weak fuzzy domain]
\label{def:weak-fuzzy-domain-enriched}
A \textbf{weak fuzzy domain} is a small \(\mathbb V\)-category \(X\) equipped with
a chosen \(\mathbb V\)-functor
\[
    \alpha:T_{\mathbb V}X\to X.
\]
Using
\[
    T_{\mathbb V}X=\operatorname{Idl}_{\mathbb V}(X),
\]
we regard \(\alpha\) as assigning to each fuzzy ideal
\[
    I:X^{\op}\to\underline{\mathbb V}
\]
a chosen candidate fuzzy ideal-join
\[
    \alpha(I)\in X.
\]
\end{definition}

\begin{definition}[Fuzzy domain]
\label{def:fuzzy-domain-enriched}
A \textbf{fuzzy domain} is a weak fuzzy domain whose structure map
\[
    \alpha:T_{\mathbb V}X\to X
\]
is a strict algebra for the fuzzy ideal monad:
\[
    \alpha\eta_X=1_X,
    \qquad
    \alpha\mu_X=\alpha T_{\mathbb V}(\alpha).
\]
\end{definition}

\begin{remark}
Weak fuzzy domains carry chosen candidate ideal-joins.  Strict fuzzy domains are
those for which these choices are coherent with principal ideals and iterated
fuzzy ideals.
\end{remark}

\subsection{The posetal specialization}
\label{subsec:weak-fuzzy-domains-posetal}

\begin{definition}[Weak fuzzy domain over \(\Lambda\)]
\label{def:weak-fuzzy-domain-posetal}
A \textbf{weak fuzzy domain over \(\Lambda\)} is a fuzzy preorder \(X\) equipped
with a \(\Lambda\)-functor
\[
    \alpha:\operatorname{Idl}_\Lambda(X)\to X.
\]
\end{definition}

\begin{definition}[Fuzzy domain over \(\Lambda\)]
\label{def:fuzzy-domain-posetal}
A \textbf{fuzzy domain over \(\Lambda\)} is a weak fuzzy domain satisfying
\[
    \alpha(\eta_X x)=x
\]
and
\[
    \alpha(\mu_X\mathcal I)=\alpha(T_\Lambda\alpha(\mathcal I)).
\]
Equivalently,
\[
    \alpha(X(-,x))=x
\]
and
\[
\alpha\left(
    x\longmapsto
    \bigjoin_{I\in\operatorname{Idl}_\Lambda(X)}
    \mathcal I(I)\tensor I(x)
\right)
=
\alpha\left(
    x\longmapsto
    \bigjoin_{I\in\operatorname{Idl}_\Lambda(X)}
    \mathcal I(I)\tensor X(x,\alpha(I))
\right).
\]
\end{definition}

\begin{theorem}[Crisp algebras are dcpos]
\label{thm:crisp-algebras-are-dcpos}
For \(\Lambda=2\), strict algebras for the ideal monad on separated preorders are
precisely dcpos, where dcpos are understood to have suprema of nonempty directed
subsets.
\end{theorem}

\begin{proof}
Let
\[
    \alpha:\operatorname{Idl}(X)\to X
\]
be a strict algebra.  The unit law says
\[
    \alpha(\downarrow x)=x.
\]
If \(I\in\operatorname{Idl}(X)\) and \(x\in I\), then
\[
    \downarrow x\subseteq I,
\]
so by monotonicity
\[
    x=\alpha(\downarrow x)\leq \alpha(I).
\]
Thus \(\alpha(I)\) is an upper bound of \(I\).  If \(y\) is any upper bound of
\(I\), then
\[
    I\subseteq\downarrow y,
\]
and again by monotonicity,
\[
    \alpha(I)\leq\alpha(\downarrow y)=y.
\]
Hence \(\alpha(I)=\bigjoin I\).  Every nonempty directed subset generates an ideal
with the same supremum, so \(X\) is a dcpo when separated.

Conversely, if \(X\) is a dcpo, define
\[
    \alpha(I)=\bigjoin I.
\]
The unit law is immediate, and the multiplication law says that the supremum of a
directed family of directed suprema is the supremum of its union.  Thus \(X\) is a
strict algebra.
\end{proof}

\begin{example}[G\"odel, product, and Lawvere weak fuzzy domains]
\label{ex:weak-fuzzy-domain-standard}
A G\"odel-valued weak fuzzy domain assigns candidate joins to finite-flat lower
weights governed by the minimum tensor.  A product-valued weak fuzzy domain assigns
candidate joins to finite-flat lower weights with multiplicative membership.  For
the Lawvere quantale, a weak fuzzy domain assigns colimits to finite-flat Lawvere
weights, giving a metric-enriched domain-like completion structure.
\end{example}

\section{The category of fuzzy domains}
\label{sec:category-fuzzy-domains}

We now record the categorical structure carried by fuzzy domains.  The point is
not to add extra structure to a fuzzy domain, but to observe that the
finite-flat doctrine already used to define fuzzy ideals also determines the
colimits, morphisms, tensor products, and internal homs of fuzzy domains.

Throughout this section, finite-flatness is understood relative to the
finite-limit doctrine fixed at the beginning of the chapter.  Thus, in each
predicate category \([X,\underline{\mathbb V}]\), the specified finite limits are
computed pointwise.  In the posetal scalar-frame specialization, this means
preservation of the terminal predicate and binary meets.

\begin{definition}[The category of fuzzy domains]
\label{def:category-fuzzy-domains}
A morphism of fuzzy domains
\[
    f:(X,\alpha_X)\longrightarrow (Y,\alpha_Y)
\]
is a \(\mathbb V\)-functor
\[
    f:X\to Y
\]
such that
\[
    f\alpha_X=\alpha_YT_{\mathbb V}(f).
\]
The ordinary category of fuzzy domains and fuzzy-domain morphisms is denoted
\[
    \mathbf{FDom}_{\mathbb V}.
\]
\end{definition}

Thus \(\mathbf{FDom}_{\mathbb V}\) is the Eilenberg--Moore category of strict
algebras for the fuzzy ideal monad.  The rest of this section identifies these
strict algebras as finite-flat-cocomplete \(\mathbb V\)-categories and then applies
Kelly's tensor-product theorem for saturated commutative classes of weights.

\subsection{Finite-flat colimits in fuzzy domains}
\label{subsec:finite-flat-colimits-fdom}

\begin{lemma}[Representables are finite-flat]
\label{lem:representables-finite-flat-fdom}
Let \(X\) be a small \(\mathbb V\)-category.  For every \(x\in X\), the
representable weight
\[
    X(-,x):X^{\op}\to\underline{\mathbb V}
\]
is finite-flat.
\end{lemma}

\begin{proof}
The colimit functional associated to \(X(-,x)\) is
\[
    \chi_{X(-,x)}(U)
    =
    \int^{y\in X}X(y,x)\tensor U(y).
\]
By enriched co-Yoneda,
\[
    \chi_{X(-,x)}(U)\cong U(x).
\]
Thus \(\chi_{X(-,x)}\) is evaluation at \(x\).  Since the specified finite
limits in \([X,\underline{\mathbb V}]\) are computed pointwise, evaluation at
\(x\) preserves them.  Hence \(X(-,x)\) is finite-flat.
\end{proof}

\begin{lemma}[Finite-flat direct images]
\label{lem:finite-flat-direct-images-fdom}
Let
\[
    D:K\to X
\]
be a \(\mathbb V\)-functor, and let
\[
    W:K^{\op}\to\underline{\mathbb V}
\]
be finite-flat.  Define the lower direct image
\[
    D_!W:X^{\op}\to\underline{\mathbb V}
\]
by
\[
    (D_!W)(x)
    =
    \int^{k\in K}W(k)\tensor X(x,Dk).
\]
Then \(D_!W\) is finite-flat.
\end{lemma}

\begin{proof}
Let
\[
    U:X\to\underline{\mathbb V}
\]
be an upper predicate on \(X\).  Then
\[
\begin{aligned}
    \chi_{D_!W}(U)
    &=
    \int^{x\in X}(D_!W)(x)\tensor U(x)
    \\
    &=
    \int^{x\in X}
        \left(
            \int^{k\in K}W(k)\tensor X(x,Dk)
        \right)\tensor U(x)
    \\
    &\cong
    \int^{k\in K}
        W(k)\tensor
        \left(
            \int^{x\in X}X(x,Dk)\tensor U(x)
        \right)
    \\
    &\cong
    \int^{k\in K}W(k)\tensor U(Dk)
    \\
    &=
    \chi_W(UD).
\end{aligned}
\]
The first isomorphism is Fubini for coends, and the second is enriched
co-Yoneda.

Precomposition
\[
    D^*:[X,\underline{\mathbb V}]\to[K,\underline{\mathbb V}],
    \qquad
    U\longmapsto UD,
\]
preserves the specified finite limits because they are computed pointwise.
Since \(W\) is finite-flat, \(\chi_W\) preserves those finite limits.  Hence
\[
    \chi_{D_!W}\cong \chi_WD^*
\]
preserves them, so \(D_!W\) is finite-flat.
\end{proof}

\begin{proposition}[Finite-flat saturation]
\label{prop:finite-flat-saturation-fdom}
Finite-flat weights are closed under finite-flat weighted colimits in presheaf
categories.

Explicitly, let
\[
    W:K^{\op}\to\underline{\mathbb V}
\]
be finite-flat, and let
\[
    H:K\to\operatorname{Idl}_{\mathbb V}(X)
\]
be a \(\mathbb V\)-functor.  Define
\[
    J:X^{\op}\to\underline{\mathbb V}
\]
pointwise by
\[
    J(x)
    =
    \int^{k\in K}W(k)\tensor Hk(x).
\]
Then \(J\) is finite-flat.
\end{proposition}

\begin{proof}
For every upper predicate
\[
    U:X\to\underline{\mathbb V},
\]
we have
\[
\begin{aligned}
    \chi_J(U)
    &=
    \int^{x\in X}J(x)\tensor U(x)
    \\
    &=
    \int^{x\in X}
        \left(
            \int^{k\in K}W(k)\tensor Hk(x)
        \right)\tensor U(x)
    \\
    &\cong
    \int^{k\in K}
        W(k)\tensor
        \left(
            \int^{x\in X}Hk(x)\tensor U(x)
        \right)
    \\
    &=
    \int^{k\in K}W(k)\tensor \chi_{Hk}(U).
\end{aligned}
\]
Thus
\[
    \chi_J(U)\cong\chi_W(\overline H(U)),
\]
where
\[
    \overline H:[X,\underline{\mathbb V}]\to[K,\underline{\mathbb V}]
\]
is defined by
\[
    \overline H(U)(k)=\chi_{Hk}(U).
\]

We first check that \(\overline H(U)\) is an upper predicate on \(K\).  Since
\(H\) is a \(\mathbb V\)-functor, for every \(k,k'\in K\) there is a morphism
\[
    K(k,k')\longrightarrow
    \operatorname{Idl}_{\mathbb V}(X)(Hk,Hk').
\]
Functoriality of the coend construction gives, naturally in \(U\),
\[
    \operatorname{Idl}_{\mathbb V}(X)(Hk,Hk')
    \longrightarrow
    \underline{\mathbb V}(\chi_{Hk}(U),\chi_{Hk'}(U)).
\]
Composing these maps gives
\[
    K(k,k')\longrightarrow
    \underline{\mathbb V}(\overline H(U)(k),\overline H(U)(k')),
\]
so \(\overline H(U)\) is indeed a \(\mathbb V\)-functor
\[
    K\to\underline{\mathbb V}.
\]

Now let \((U_i)\) be a finite diagram of upper predicates whose specified limit
exists in the finite-limit doctrine.  Since each \(Hk\) is finite-flat,
\[
    \chi_{Hk}(\lim_i U_i)
    \cong
    \lim_i\chi_{Hk}(U_i).
\]
Since specified finite limits in \([K,\underline{\mathbb V}]\) are computed
pointwise, this gives
\[
    \overline H(\lim_i U_i)\cong \lim_i\overline H(U_i).
\]
Thus \(\overline H\) preserves the specified finite limits.  Since \(W\) is
finite-flat, \(\chi_W\) preserves them.  Hence
\[
    \chi_J\cong\chi_W\overline H
\]
preserves the specified finite limits, so \(J\) is finite-flat.
\end{proof}

\begin{corollary}[Flattening is finite-flat]
\label{cor:flattening-finite-flat-fdom}
If
\[
    \mathcal I:
    \operatorname{Idl}_{\mathbb V}(X)^{\op}\to\underline{\mathbb V}
\]
is finite-flat, then the flattened weight
\[
    \mu_X(\mathcal I)(x)
    =
    \int^{I\in\operatorname{Idl}_{\mathbb V}(X)}
        \mathcal I(I)\tensor I(x)
\]
is finite-flat.
\end{corollary}

\begin{proof}
Apply Proposition~\ref{prop:finite-flat-saturation-fdom} to the identity diagram
\[
    1_{\operatorname{Idl}_{\mathbb V}(X)}:
    \operatorname{Idl}_{\mathbb V}(X)\to\operatorname{Idl}_{\mathbb V}(X).
\]
\end{proof}

\begin{proposition}[The KZ comparison for the ideal monad]
\label{prop:ideal-monad-kz-fdom}
For every small \(\mathbb V\)-category \(X\), there is an enriched natural
transformation
\[
    T_{\mathbb V}\eta_X\Longrightarrow \eta_{T_{\mathbb V}X}.
\]
Equivalently, for every finite-flat weight
\[
    I:X^{\op}\to\underline{\mathbb V},
\]
there is a canonical morphism of weights
\[
    (T_{\mathbb V}\eta_X)(I)
    \longrightarrow
    \operatorname{Idl}_{\mathbb V}(X)(-,I).
\]
Together with the monad unit and multiplication, these comparisons make the ideal
monad a KZ, or lax-idempotent, monad.
\end{proposition}

\begin{proof}
Let
\[
    I,J\in\operatorname{Idl}_{\mathbb V}(X).
\]
The value of \((T_{\mathbb V}\eta_X)(I)\) at \(J\) is
\[
    (T_{\mathbb V}\eta_X)(I)(J)
    =
    \int^{x\in X}
        I(x)\tensor
        \operatorname{Idl}_{\mathbb V}(X)(J,X(-,x)).
\]
By enriched Yoneda,
\[
    I(x)
    \cong
    [X^{\op},\underline{\mathbb V}](X(-,x),I).
\]
Composition in the presheaf category gives morphisms
\[
    \operatorname{Idl}_{\mathbb V}(X)(J,X(-,x))
    \tensor
    [X^{\op},\underline{\mathbb V}](X(-,x),I)
    \longrightarrow
    \operatorname{Idl}_{\mathbb V}(X)(J,I).
\]
Tensoring with the Yoneda identification and integrating over \(x\) gives
\[
    (T_{\mathbb V}\eta_X)(I)(J)
    \longrightarrow
    \operatorname{Idl}_{\mathbb V}(X)(J,I).
\]
This is natural in \(J\), and therefore defines a morphism of weights
\[
    (T_{\mathbb V}\eta_X)(I)
    \longrightarrow
    \operatorname{Idl}_{\mathbb V}(X)(-,I).
\]
It is also enriched-natural in \(I\), since it is induced by composition in
\([X^{\op},\underline{\mathbb V}]\).

We now record the KZ consequences explicitly.  The monad laws give
\[
    \mu_X\,T_{\mathbb V}\eta_X=1_{T_{\mathbb V}X},
    \qquad
    \mu_X\,\eta_{T_{\mathbb V}X}=1_{T_{\mathbb V}X}.
\]
The comparison just constructed gives
\[
    T_{\mathbb V}\eta_X\Longrightarrow \eta_{T_{\mathbb V}X}.
\]
Composing with \(\mu_X\) and using the first monad law gives the identity on
\(T_{\mathbb V}X\).  Equivalently, the comparison exhibits
\[
    \mu_X\dashv \eta_{T_{\mathbb V}X}
\]
in the enriched order of endomorphisms.  The remaining triangle inequality is the
statement that every finite-flat ideal of finite-flat ideals is contained in the
principal ideal generated by its flattening:
\[
    1_{T_{\mathbb V}T_{\mathbb V}X}
    \Longrightarrow
    \eta_{T_{\mathbb V}X}\mu_X.
\]
On a finite-flat weight
\[
    \mathcal I:
    \operatorname{Idl}_{\mathbb V}(X)^{\op}\to\underline{\mathbb V},
\]
this is the morphism
\[
    \mathcal I(J)
    \longrightarrow
    \operatorname{Idl}_{\mathbb V}(X)(J,\mu_X\mathcal I)
\]
obtained by composing
\[
    \mathcal I(J)\tensor J(x)
    \longrightarrow
    \int^{I}\mathcal I(I)\tensor I(x)
    =
    \mu_X(\mathcal I)(x)
\]
and transposing in the presheaf hom.  Naturality and the triangle identities are
the associativity and unitality of composition in
\([X^{\op},\underline{\mathbb V}]\), together with Fubini for the coends defining
flattening.  Hence the ideal monad is KZ.
\end{proof}

\begin{theorem}[The algebra map is left adjoint to Yoneda]
\label{thm:algebra-map-left-adjoint-y-fdom}
Let
\[
    \alpha:\operatorname{Idl}_{\mathbb V}(X)\to X
\]
be a strict algebra structure for the fuzzy ideal monad.  Then
\[
    \alpha\dashv \eta_X
\]
with identity counit
\[
    \alpha\eta_X=1_X.
\]
Consequently, enriched-naturally in
\(I\in\operatorname{Idl}_{\mathbb V}(X)\) and \(x\in X\), there is an isomorphism
\[
    X(\alpha I,x)
    \cong
    \operatorname{Idl}_{\mathbb V}(X)(I,X(-,x)).
\]
\end{theorem}

\begin{proof}
The algebra unit law gives
\[
    \alpha\eta_X=1_X,
\]
so the counit of the desired adjunction is the identity.

It remains to construct the unit
\[
    1_{\operatorname{Idl}_{\mathbb V}(X)}
    \longrightarrow
    \eta_X\alpha.
\]
By Proposition~\ref{prop:ideal-monad-kz-fdom}, we have
\[
    T_{\mathbb V}\eta_X\Longrightarrow \eta_{T_{\mathbb V}X}.
\]
Applying \(T_{\mathbb V}\alpha\) gives
\[
    T_{\mathbb V}\alpha\;T_{\mathbb V}\eta_X
    \longrightarrow
    T_{\mathbb V}\alpha\;\eta_{T_{\mathbb V}X}.
\]
The left-hand side is
\[
    T_{\mathbb V}(\alpha\eta_X)=T_{\mathbb V}(1_X)=1_{\operatorname{Idl}_{\mathbb V}(X)}.
\]
By naturality of \(\eta\), the right-hand side is
\[
    \eta_X\alpha.
\]
Thus we obtain
\[
    1_{\operatorname{Idl}_{\mathbb V}(X)}
    \longrightarrow
    \eta_X\alpha.
\]
The KZ coherence supplies the triangle identities.  Therefore
\[
    \alpha\dashv\eta_X.
\]
The displayed hom-isomorphism is the enriched hom-form of this adjunction.
\end{proof}

\begin{theorem}[Fuzzy domains and finite-flat colimits]
\label{thm:fuzzy-domains-ff-colimits}
Let \(X\) be a small \(\mathbb V\)-category.  To give \(X\) the structure of a
fuzzy domain is equivalently to give a strict coherent choice of all finite-flat
weighted colimits in \(X\).

Under this equivalence, the algebra map
\[
    \alpha_X:\operatorname{Idl}_{\mathbb V}(X)\to X
\]
sends a finite-flat weight \(I:X^{\op}\to\underline{\mathbb V}\) to the
\(I\)-weighted colimit of the identity diagram:
\[
    \alpha_X(I)=I*1_X.
\]
More generally, if
\[
    D:K\to X
\]
is a diagram and
\[
    W:K^{\op}\to\underline{\mathbb V}
\]
is finite-flat, then
\[
    W*D=\alpha_X(D_!W),
\]
where
\[
    (D_!W)(x)
    =
    \int^{k\in K}W(k)\tensor X(x,Dk).
\]

A \(\mathbb V\)-functor between fuzzy domains is a fuzzy-domain morphism if and
only if it preserves finite-flat weighted colimits.
\end{theorem}

\begin{proof}
Suppose first that
\[
    \alpha_X:\operatorname{Idl}_{\mathbb V}(X)\to X
\]
is a strict algebra.  By
Theorem~\ref{thm:algebra-map-left-adjoint-y-fdom},
\[
    \alpha_X\dashv\eta_X.
\]
Hence, for every finite-flat weight \(I\) and every \(x\in X\),
\[
\begin{aligned}
    X(\alpha_XI,x)
    &\cong
    \operatorname{Idl}_{\mathbb V}(X)(I,\eta_Xx)
    \\
    &=
    [X^{\op},\underline{\mathbb V}](I,X(-,x)).
\end{aligned}
\]
This is the universal property of the \(I\)-weighted colimit of the identity
diagram.  Thus
\[
    \alpha_X(I)=I*1_X.
\]

Now let \(D:K\to X\), and let \(W:K^{\op}\to\underline{\mathbb V}\) be
finite-flat.  By Lemma~\ref{lem:finite-flat-direct-images-fdom}, \(D_!W\) is
finite-flat.  For every \(x\in X\),
\[
\begin{aligned}
    X(\alpha_X(D_!W),x)
    &\cong
    [X^{\op},\underline{\mathbb V}](D_!W,X(-,x))
    \\
    &\cong
    [K^{\op},\underline{\mathbb V}](W,X(D-,x)).
\end{aligned}
\]
The second isomorphism is the weighted left-Kan-extension adjunction defining
\(D_!W\).  This is the universal property of the weighted colimit \(W*D\).
Therefore
\[
    W*D=\alpha_X(D_!W)
\]
under the strictification convention.

Conversely, suppose \(X\) is equipped with a strict coherent choice of all
finite-flat weighted colimits.  Define
\[
    \alpha_X(I)=I*1_X.
\]
The representing isomorphisms
\[
    X(\alpha_XI,x)
    \cong
    [X^{\op},\underline{\mathbb V}](I,X(-,x))
\]
are enriched-natural in \(I\) and \(x\).  Hence \(I\mapsto \alpha_XI\) extends
uniquely to a \(\mathbb V\)-functor
\[
    \alpha_X:\operatorname{Idl}_{\mathbb V}(X)\to X
\]
left adjoint to the Yoneda embedding.

For a representable weight,
\[
    X(-,x)*1_X=x,
\]
so
\[
    \alpha_X\eta_X=1_X.
\]
For a finite-flat weight
\[
    \mathcal I:
    \operatorname{Idl}_{\mathbb V}(X)^{\op}\to\underline{\mathbb V},
\]
the left-hand side of the multiplication law is
\[
    \alpha_X\mu_X(\mathcal I)=\mu_X(\mathcal I)*1_X.
\]
The right-hand side is
\[
    \alpha_XT_{\mathbb V}(\alpha_X)(\mathcal I)
    =
    (T_{\mathbb V}\alpha_X)(\mathcal I)*1_X.
\]
The weight \(T_{\mathbb V}(\alpha_X)(\mathcal I)\) is the lower direct image of
\(\mathcal I\) along
\[
    \alpha_X:\operatorname{Idl}_{\mathbb V}(X)\to X.
\]
Hence its colimit against \(1_X\) is the same weighted colimit as
\(\mathcal I\) against the diagram \(\alpha_X\):
\[
    (T_{\mathbb V}\alpha_X)(\mathcal I)*1_X
    \cong
    \mathcal I*\alpha_X.
\]
On the other hand,
\[
    \mu_X(\mathcal I)(x)
    =
    \int^{I}\mathcal I(I)\tensor I(x)
\]
is the flattened weight computing the same iterated finite-flat colimit.
Strict coherence of the chosen finite-flat colimits gives
\[
    \mu_X(\mathcal I)*1_X
    =
    \mathcal I*(I\mapsto I*1_X)
    =
    \mathcal I*\alpha_X.
\]
Therefore
\[
    \alpha_X\mu_X=\alpha_XT_{\mathbb V}(\alpha_X),
\]
so \(\alpha_X\) is a strict algebra.

Finally, let
\[
    f:(X,\alpha_X)\to(Y,\alpha_Y)
\]
be a \(\mathbb V\)-functor.  If \(f\) is a fuzzy-domain morphism, then for every
finite-flat \(I:X^{\op}\to\underline{\mathbb V}\),
\[
    f(I*1_X)
    =
    f\alpha_X(I)
    =
    \alpha_YT_{\mathbb V}(f)(I)
    =
    T_{\mathbb V}(f)(I)*1_Y.
\]
Thus \(f\) preserves finite-flat colimits of identity diagrams, and hence, by
the direct-image formula, all finite-flat weighted colimits.

Conversely, if \(f\) preserves finite-flat weighted colimits, then for every
finite-flat \(I\),
\[
    f\alpha_X(I)
    =
    f(I*1_X)
    =
    T_{\mathbb V}(f)(I)*1_Y
    =
    \alpha_YT_{\mathbb V}(f)(I).
\]
Therefore
\[
    f\alpha_X=\alpha_YT_{\mathbb V}(f),
\]
so \(f\) is a fuzzy-domain morphism.
\end{proof}

\subsection{Free fuzzy domains and limits}
\label{subsec:free-fuzzy-domains-limits}

\begin{theorem}[Free fuzzy domains]
\label{thm:free-fuzzy-domains}
The forgetful functor
\[
    U:\mathbf{FDom}_{\mathbb V}\to \mathbb V\text{-}\mathbf{Cat}
\]
has a left adjoint
\[
    F:\mathbb V\text{-}\mathbf{Cat}\to\mathbf{FDom}_{\mathbb V},
    \qquad
    F(X)=\operatorname{Idl}_{\mathbb V}(X),
\]
where \(F(X)\) has algebra structure
\[
    \mu_X:
    \operatorname{Idl}_{\mathbb V}(\operatorname{Idl}_{\mathbb V}(X))
    \to
    \operatorname{Idl}_{\mathbb V}(X).
\]
Thus, for every small \(\mathbb V\)-category \(X\) and every fuzzy domain \(A\),
there is a natural bijection
\[
    \mathbf{FDom}_{\mathbb V}(\operatorname{Idl}_{\mathbb V}(X),A)
    \cong
    \mathbb V\text{-}\mathbf{Cat}(X,UA).
\]
\end{theorem}

\begin{proof}
This is the free-algebra adjunction of the monad \(T_{\mathbb V}\).  Explicitly,
if
\[
    f:X\to UA
\]
is a \(\mathbb V\)-functor and
\[
    \alpha_A:\operatorname{Idl}_{\mathbb V}(A)\to A
\]
is the algebra map of \(A\), define
\[
    \overline f:
    \operatorname{Idl}_{\mathbb V}(X)\to A
\]
by
\[
    \overline f(I)
    =
    \alpha_A(T_{\mathbb V}f(I)).
\]
Equivalently,
\[
    \overline f(I)=I*f.
\]
By Theorem~\ref{thm:fuzzy-domains-ff-colimits}, \(\overline f\) preserves
finite-flat weighted colimits, so it is a fuzzy-domain morphism.  Moreover,
\[
    \overline f(\eta_Xx)
    =
    \alpha_A(\eta_A(fx))
    =
    fx.
\]
Thus
\[
    \overline f\eta_X=f.
\]

Conversely, if
\[
    g:\operatorname{Idl}_{\mathbb V}(X)\to A
\]
is a fuzzy-domain morphism, then \(g\) preserves finite-flat weighted colimits.
For every finite-flat weight \(I:X^{\op}\to\underline{\mathbb V}\), the object
\(I\in\operatorname{Idl}_{\mathbb V}(X)\) is the \(I\)-weighted colimit of the
Yoneda diagram:
\[
    I\cong I*\eta_X.
\]
Therefore
\[
    g(I)
    \cong
    g(I*\eta_X)
    \cong
    I*(g\eta_X).
\]
Hence
\[
    g=\overline{g\eta_X}
\]
under the strictification convention.  The two constructions are inverse and
natural.
\end{proof}

\begin{corollary}[Monadicity]
\label{cor:fdom-monadicity}
The forgetful functor
\[
    U:\mathbf{FDom}_{\mathbb V}\to \mathbb V\text{-}\mathbf{Cat}
\]
is monadic, and the induced monad is the fuzzy ideal monad
\[
    T_{\mathbb V}.
\]
Thus
\[
    \mathbf{FDom}_{\mathbb V}=(\mathbb V\text{-}\mathbf{Cat})^{T_{\mathbb V}}.
\]
\end{corollary}

\begin{proof}
By definition, \(\mathbf{FDom}_{\mathbb V}\) is the Eilenberg--Moore category of
strict algebras for \(T_{\mathbb V}\).  The adjunction of
Theorem~\ref{thm:free-fuzzy-domains} has underlying monad \(T_{\mathbb V}\), with
Yoneda unit and flattening multiplication.
\end{proof}

\begin{theorem}[Small conical limits]
\label{thm:fdom-limits}
The category \(\mathbf{FDom}_{\mathbb V}\) has all small conical limits whose
underlying limits exist in
\[
    \mathbb V\text{-}\mathbf{Cat}.
\]
These limits are created by the forgetful functor
\[
    U:\mathbf{FDom}_{\mathbb V}\to \mathbb V\text{-}\mathbf{Cat}.
\]
\end{theorem}

\begin{proof}
Eilenberg--Moore categories have all limits that exist in the base category, and
the forgetful functor creates them.

Explicitly, let
\[
    D:J\to\mathbf{FDom}_{\mathbb V}
\]
be a small diagram, and suppose that
\[
    L=\lim_{j\in J}UD_j
\]
exists in \(\mathbb V\text{-}\mathbf{Cat}\), with projections
\[
    p_j:L\to UD_j.
\]
There is a unique algebra structure
\[
    \alpha_L:T_{\mathbb V}L\to L
\]
such that
\[
    p_j\alpha_L=\alpha_jT_{\mathbb V}(p_j)
\]
for every \(j\).  With this algebra structure, \(L\) satisfies the universal
property of the limit in \(\mathbf{FDom}_{\mathbb V}\).
\end{proof}

\subsection{Tensor products and internal homs}
\label{subsec:tensor-products-internal-homs-fdom}

\begin{definition}[External product of weights]
\label{def:external-product-weights-fdom}
Let
\[
    W:K^{\op}\to\underline{\mathbb V},
    \qquad
    Q:J^{\op}\to\underline{\mathbb V}
\]
be weights.  Their external product is the weight
\[
    W\boxtimes Q:(K\tensor J)^{\op}\to\underline{\mathbb V}
\]
defined by
\[
    (W\boxtimes Q)(k,j)=W(k)\tensor Q(j).
\]
\end{definition}

\begin{proposition}[External products of finite-flat weights]
\label{prop:external-products-finite-flat-fdom}
If \(W\) and \(Q\) are finite-flat, then
\[
    W\boxtimes Q
\]
is finite-flat.
\end{proposition}

\begin{proof}
Let
\[
    E:K\tensor J\to\underline{\mathbb V}
\]
be an upper predicate.  Then
\[
\begin{aligned}
    \chi_{W\boxtimes Q}(E)
    &=
    \int^{(k,j)\in K\tensor J}
        W(k)\tensor Q(j)\tensor E(k,j)
    \\
    &\cong
    \int^{k\in K}
        W(k)\tensor
        \left(
            \int^{j\in J}Q(j)\tensor E(k,j)
        \right).
\end{aligned}
\]
For each \(k\), the inner expression is
\[
    \chi_Q(E(k,-)).
\]

Since \(E:K\tensor J\to\underline{\mathbb V}\) is a \(\mathbb V\)-functor, the
assignment
\[
    k\longmapsto E(k,-)
\]
defines a \(\mathbb V\)-functor
\[
    K\to[J,\underline{\mathbb V}].
\]
Composing with
\[
    \chi_Q:[J,\underline{\mathbb V}]\to\underline{\mathbb V}
\]
gives an upper predicate
\[
    k\longmapsto \chi_Q(E(k,-))
\]
on \(K\).  Since \(Q\) is finite-flat, this construction preserves the specified
finite limits in the \(J\)-variable.  Since \(W\) is finite-flat, applying
\(\chi_W\) preserves the specified finite limits in the remaining \(K\)-variable.
Hence \(\chi_{W\boxtimes Q}\) preserves the specified finite limits, so
\(W\boxtimes Q\) is finite-flat.
\end{proof}

\begin{proposition}[Finite-flat Fubini]
\label{prop:finite-flat-fubini-fdom}
Let
\[
    W:K^{\op}\to\underline{\mathbb V},
    \qquad
    Q:J^{\op}\to\underline{\mathbb V}
\]
be finite-flat.  For every diagram
\[
    E:K\tensor J\to C
\]
in a fuzzy domain \(C\), there are canonical isomorphisms
\[
    W*\bigl(k\mapsto Q*(j\mapsto E(k,j))\bigr)
    \cong
    (W\boxtimes Q)*E
    \cong
    Q*\bigl(j\mapsto W*(k\mapsto E(k,j))\bigr).
\]
\end{proposition}

\begin{proof}
These are the Fubini isomorphisms for enriched weighted colimits.  By
Proposition~\ref{prop:external-products-finite-flat-fdom}, the external product
\(W\boxtimes Q\) is finite-flat.  Since \(C\) is a fuzzy domain, all three
finite-flat weighted colimits appearing above exist.
\end{proof}

\begin{definition}[Internal hom fuzzy domain]
\label{def:internal-hom-fdom}
For fuzzy domains \(B\) and \(C\), define
\[
    [B,C]_{\mathrm{ff}}
\]
to be the full \(\mathbb V\)-subcategory of the enriched functor category
\([B,C]\) spanned by the finite-flat-colimit-preserving \(\mathbb V\)-functors
\[
    B\to C.
\]
\end{definition}

\begin{proposition}[Internal homs are fuzzy domains]
\label{prop:internal-homs-fdom}
For fuzzy domains \(B\) and \(C\), the \(\mathbb V\)-category
\[
    [B,C]_{\mathrm{ff}}
\]
is a fuzzy domain.  Its finite-flat weighted colimits are computed pointwise in
\(C\).
\end{proposition}

\begin{proof}
Let
\[
    W:K^{\op}\to\underline{\mathbb V}
\]
be finite-flat, and let
\[
    H:K\to[B,C]_{\mathrm{ff}}
\]
be a diagram.  Since \(C\) is a fuzzy domain, the pointwise formula
\[
    (W*H)(b)=W*(k\mapsto Hk(b))
\]
defines the \(W\)-weighted colimit of \(H\) in the ambient functor category
\([B,C]\).

It remains to show that \(W*H\) preserves finite-flat weighted colimits in
\(B\).  Let
\[
    Q:J^{\op}\to\underline{\mathbb V}
\]
be finite-flat, and let
\[
    D:J\to B
\]
be a diagram.  Then
\[
\begin{aligned}
    (W*H)(Q*D)
    &\cong
    W*\bigl(k\mapsto Hk(Q*D)\bigr)
    \\
    &\cong
    W*\bigl(k\mapsto Q*(j\mapsto Hk(Dj))\bigr)
    \\
    &\cong
    Q*\bigl(j\mapsto W*(k\mapsto Hk(Dj))\bigr)
    \\
    &\cong
    Q*((W*H)D).
\end{aligned}
\]
The second isomorphism uses that each \(Hk:B\to C\) preserves finite-flat
weighted colimits.  The third isomorphism is
Proposition~\ref{prop:finite-flat-fubini-fdom}.  Hence \(W*H\) preserves
finite-flat weighted colimits in \(B\).

Therefore \([B,C]_{\mathrm{ff}}\) is closed under finite-flat weighted colimits
inside \([B,C]\).  Since \([B,C]_{\mathrm{ff}}\) is full in \([B,C]\), the
ambient weighted-colimit universal property restricts to
\([B,C]_{\mathrm{ff}}\).  By
Theorem~\ref{thm:fuzzy-domains-ff-colimits},
\([B,C]_{\mathrm{ff}}\) is a fuzzy domain.
\end{proof}

\begin{definition}[Finite-flat bimorphism]
\label{def:finite-flat-bimorphism-fdom}
Let \(A,B,C\) be fuzzy domains.  A finite-flat bimorphism
\[
    M:A\tensor B\to C
\]
is a \(\mathbb V\)-functor such that:
\begin{enumerate}
    \item for each \(a\in A\), the partial functor
    \[
        M(a,-):B\to C
    \]
    preserves finite-flat weighted colimits;
    \item for each \(b\in B\), the partial functor
    \[
        M(-,b):A\to C
    \]
    preserves finite-flat weighted colimits.
\end{enumerate}
We write
\[
    \operatorname{Bimor}_{\mathrm{ff}}(A,B;C)
\]
for the set of finite-flat bimorphisms from \(A,B\) to \(C\).
\end{definition}

\begin{proposition}[Currying finite-flat bimorphisms]
\label{prop:currying-bimorphisms-fdom}
For fuzzy domains \(A,B,C\), there is a natural bijection
\[
    \operatorname{Bimor}_{\mathrm{ff}}(A,B;C)
    \cong
    \mathbf{FDom}_{\mathbb V}(A,[B,C]_{\mathrm{ff}}).
\]
\end{proposition}

\begin{proof}
Let
\[
    M:A\tensor B\to C
\]
be a finite-flat bimorphism.  By enriched currying, \(M\) corresponds to a
\(\mathbb V\)-functor
\[
    \widehat M:A\to[B,C],
    \qquad
    \widehat M(a)(b)=M(a,b).
\]
Since \(M\) preserves finite-flat weighted colimits in the \(B\)-variable, each
\[
    \widehat M(a):B\to C
\]
lies in \([B,C]_{\mathrm{ff}}\).  Thus \(\widehat M\) factors through
\[
    [B,C]_{\mathrm{ff}}\subseteq[B,C].
\]
Since \(M\) preserves finite-flat weighted colimits in the \(A\)-variable and
finite-flat colimits in \([B,C]_{\mathrm{ff}}\) are computed pointwise,
\[
    \widehat M:A\to[B,C]_{\mathrm{ff}}
\]
preserves finite-flat weighted colimits.  Hence \(\widehat M\) is a
fuzzy-domain morphism.

Conversely, let
\[
    H:A\to[B,C]_{\mathrm{ff}}
\]
be a fuzzy-domain morphism.  Uncurrying gives
\[
    \widetilde H:A\tensor B\to C,
    \qquad
    \widetilde H(a,b)=H(a)(b).
\]
For each \(a\), the partial functor
\[
    \widetilde H(a,-)=H(a)
\]
preserves finite-flat weighted colimits because \(H(a)\) is an object of
\([B,C]_{\mathrm{ff}}\).  For each \(b\), the partial functor
\[
    \widetilde H(-,b):A\to C
\]
preserves finite-flat weighted colimits because \(H\) does and colimits in
\([B,C]_{\mathrm{ff}}\) are computed pointwise.  Hence \(\widetilde H\) is a
finite-flat bimorphism.

The two constructions are inverse by the enriched tensor--hom adjunction.
\end{proof}

\begin{theorem}[Kelly tensor product for fuzzy domains]
\label{thm:kelly-tensor-product-fdom}
For fuzzy domains \(A\) and \(B\), there is a fuzzy domain
\[
    A\boxtimes_{\mathrm{ff}}B
\]
and a finite-flat bimorphism
\[
    \kappa_{A,B}:A\tensor B\to A\boxtimes_{\mathrm{ff}}B
\]
such that, for every fuzzy domain \(C\), composition with \(\kappa_{A,B}\)
induces a natural bijection
\[
    \mathbf{FDom}_{\mathbb V}(A\boxtimes_{\mathrm{ff}}B,C)
    \cong
    \operatorname{Bimor}_{\mathrm{ff}}(A,B;C).
\]
\end{theorem}

\begin{proof}
By Lemma~\ref{lem:representables-finite-flat-fdom},
Lemma~\ref{lem:finite-flat-direct-images-fdom}, and
Proposition~\ref{prop:finite-flat-saturation-fdom}, the assignment
\[
    X\longmapsto \operatorname{Idl}_{\mathbb V}(X)
\]
contains the representables, is stable under lower direct image along
\(\mathbb V\)-functors, and is closed under finite-flat weighted colimits in
presheaf categories.  Hence the finite-flat weights form a saturated doctrine.

Moreover, Proposition~\ref{prop:external-products-finite-flat-fdom} and
Proposition~\ref{prop:finite-flat-fubini-fdom} show that the finite-flat
doctrine is commutative: finite-flat weighted colimits commute with finite-flat
weighted colimits.

Kelly's tensor-product theorem for categories equipped with a saturated
commutative class of specified weighted colimits therefore applies to the
finite-flat doctrine.  The representing object for separately
finite-flat-colimit-preserving \(\mathbb V\)-functors
\[
    A\tensor B\to C
\]
is denoted
\[
    A\boxtimes_{\mathrm{ff}}B.
\]
Its universal map
\[
    \kappa_{A,B}:A\tensor B\to A\boxtimes_{\mathrm{ff}}B
\]
is a finite-flat bimorphism, and the displayed bijection is exactly the
representing universal property.
\end{proof}

\begin{corollary}[Tensor product of free fuzzy domains]
\label{cor:tensor-free-fdom}
For small \(\mathbb V\)-categories \(X\) and \(Y\), there is a canonical
equivalence
\[
    \operatorname{Idl}_{\mathbb V}(X)
    \boxtimes_{\mathrm{ff}}
    \operatorname{Idl}_{\mathbb V}(Y)
    \simeq
    \operatorname{Idl}_{\mathbb V}(X\tensor Y).
\]
\end{corollary}

\begin{proof}
For every fuzzy domain \(C\), there are natural bijections
\[
\begin{aligned}
\mathbf{FDom}_{\mathbb V}
\bigl(
    \operatorname{Idl}_{\mathbb V}(X)
    \boxtimes_{\mathrm{ff}}
    \operatorname{Idl}_{\mathbb V}(Y),
    C
\bigr)
&\cong
\operatorname{Bimor}_{\mathrm{ff}}
\bigl(
    \operatorname{Idl}_{\mathbb V}(X),
    \operatorname{Idl}_{\mathbb V}(Y);
    C
\bigr)
\\
&\cong
\mathbb V\text{-}\mathbf{Cat}(X\tensor Y,UC)
\\
&\cong
\mathbf{FDom}_{\mathbb V}
\bigl(
    \operatorname{Idl}_{\mathbb V}(X\tensor Y),
    C
\bigr).
\end{aligned}
\]
The middle bijection sends a bimorphism to its restriction along
\[
    \eta_X\tensor\eta_Y:
    X\tensor Y\to
    \operatorname{Idl}_{\mathbb V}(X)\tensor\operatorname{Idl}_{\mathbb V}(Y).
\]
Conversely, a \(\mathbb V\)-functor
\[
    m:X\tensor Y\to UC
\]
extends uniquely to the finite-flat bimorphism
\[
    \widetilde m:
    \operatorname{Idl}_{\mathbb V}(X)\tensor\operatorname{Idl}_{\mathbb V}(Y)\to C
\]
defined by
\[
    \widetilde m(I,J)
    =
    I*\bigl(x\mapsto J*(y\mapsto m(x,y))\bigr).
\]
By Proposition~\ref{prop:external-products-finite-flat-fdom} and
Proposition~\ref{prop:finite-flat-fubini-fdom}, this is equivalently
\[
    \widetilde m(I,J)\cong (I\boxtimes J)*m.
\]
Thus \(\widetilde m\) preserves finite-flat weighted colimits separately in
\(I\) and \(J\), and it is the unique such extension of \(m\).  By
representability, the displayed equivalence follows.
\end{proof}

\begin{corollary}[Tensor unit]
\label{cor:tensor-unit-fdom}
Let
\[
    \mathbf 1
\]
be the terminal small \(\mathbb V\)-category.  The tensor unit for fuzzy domains is
\[
    \operatorname{Idl}_{\mathbb V}(\mathbf 1).
\]
\end{corollary}

\begin{proof}
For every fuzzy domain \(C\),
\[
\begin{aligned}
    \mathbf{FDom}_{\mathbb V}
    \bigl(
        \operatorname{Idl}_{\mathbb V}(\mathbf 1)
        \boxtimes_{\mathrm{ff}}A,
        C
    \bigr)
    &\cong
    \operatorname{Bimor}_{\mathrm{ff}}
    \bigl(
        \operatorname{Idl}_{\mathbb V}(\mathbf 1),
        A;
        C
    \bigr)
    \\
    &\cong
    \mathbf{FDom}_{\mathbb V}(A,C).
\end{aligned}
\]
The second bijection uses that \(\operatorname{Idl}_{\mathbb V}(\mathbf 1)\) is
freely generated as a fuzzy domain by the unique object of \(\mathbf 1\).  Hence
\[
    \operatorname{Idl}_{\mathbb V}(\mathbf 1)\boxtimes_{\mathrm{ff}}A
    \simeq
    A.
\]
The right unit law is analogous.
\end{proof}

\begin{theorem}[Symmetric monoidal closure of fuzzy domains]
\label{thm:fdom-smc}
The category
\[
    \mathbf{FDom}_{\mathbb V}
\]
is symmetric monoidal closed.  Its tensor product is
\[
    A\boxtimes_{\mathrm{ff}}B,
\]
its tensor unit is
\[
    \operatorname{Idl}_{\mathbb V}(\mathbf 1),
\]
and its internal hom is
\[
    [B,C]_{\mathrm{ff}}.
\]
The closed structure is the natural bijection
\[
    \mathbf{FDom}_{\mathbb V}(A\boxtimes_{\mathrm{ff}}B,C)
    \cong
    \mathbf{FDom}_{\mathbb V}(A,[B,C]_{\mathrm{ff}}).
\]
\end{theorem}

\begin{proof}
By Theorem~\ref{thm:kelly-tensor-product-fdom},
\[
    \mathbf{FDom}_{\mathbb V}(A\boxtimes_{\mathrm{ff}}B,C)
    \cong
    \operatorname{Bimor}_{\mathrm{ff}}(A,B;C).
\]
By Proposition~\ref{prop:currying-bimorphisms-fdom},
\[
    \operatorname{Bimor}_{\mathrm{ff}}(A,B;C)
    \cong
    \mathbf{FDom}_{\mathbb V}(A,[B,C]_{\mathrm{ff}}).
\]
Composing these bijections gives the closed structure.

The symmetry of \(\boxtimes_{\mathrm{ff}}\) follows from the symmetry of
\[
    A\tensor B\cong B\tensor A
\]
and from the symmetric nature of the condition of preserving finite-flat
weighted colimits separately in the two variables.

Associativity follows by representability.  For every fuzzy domain \(D\), both
\[
    (A\boxtimes_{\mathrm{ff}}B)\boxtimes_{\mathrm{ff}}C
\]
and
\[
    A\boxtimes_{\mathrm{ff}}(B\boxtimes_{\mathrm{ff}}C)
\]
represent trilinear \(\mathbb V\)-functors
\[
    A\tensor B\tensor C\to D
\]
that preserve finite-flat weighted colimits separately in each variable.  Hence
they are canonically equivalent.

The tensor unit is \(\operatorname{Idl}_{\mathbb V}(\mathbf 1)\) by
Corollary~\ref{cor:tensor-unit-fdom}.  Therefore \(\mathbf{FDom}_{\mathbb V}\) is
symmetric monoidal closed.
\end{proof}

\begin{corollary}[Evaluation]
\label{cor:evaluation-fdom}
For fuzzy domains \(B\) and \(C\), there is an evaluation morphism
\[
    [B,C]_{\mathrm{ff}}\boxtimes_{\mathrm{ff}}B
    \longrightarrow
    C.
\]
On generators it is given by
\[
    (f,b)\longmapsto f(b).
\]
\end{corollary}

\begin{proof}
The ordinary enriched evaluation
\[
    [B,C]\tensor B\to C
\]
restricts to a finite-flat bimorphism
\[
    [B,C]_{\mathrm{ff}}\tensor B\to C.
\]
It is finite-flat-colimit-preserving in the \(B\)-variable because every object
of \([B,C]_{\mathrm{ff}}\) preserves finite-flat weighted colimits.  It is
finite-flat-colimit-preserving in the \([B,C]_{\mathrm{ff}}\)-variable because
finite-flat colimits there are computed pointwise in \(C\).  By the universal
property of \(\boxtimes_{\mathrm{ff}}\), this bimorphism induces the displayed
fuzzy-domain morphism.
\end{proof}

\subsection{Posetal and crisp specializations}
\label{subsec:posetal-crisp-fdom-category}

In the posetal specialization, a fuzzy domain over \(\Lambda\) is a fuzzy
preorder \(X\) equipped with a strict algebra
\[
    \alpha_X:\operatorname{Idl}_{\Lambda}(X)\to X.
\]
The map \(\alpha_X\) sends a fuzzy ideal \(I\) to its fuzzy ideal-join.

If
\[
    D:K\to X
\]
is a \(\Lambda\)-functor and
\[
    W:K^{\op}\to\Lambda
\]
is a fuzzy ideal of \(K\), that is, a finite-flat lower fuzzy predicate on
\(K\), then
\[
    W*D=\alpha_X(D_!W),
\]
where
\[
    (D_!W)(x)
    =
    \bigjoin_{k\in K}W(k)\tensor X(x,Dk).
\]
A morphism of fuzzy domains
\[
    f:X\to Y
\]
is a \(\Lambda\)-functor satisfying
\[
    f\alpha_X=\alpha_YT_{\Lambda}(f).
\]
Equivalently, for every fuzzy ideal \(I\) of \(X\),
\[
    f(\alpha_X(I))
    =
    \alpha_Y(T_{\Lambda}f(I)).
\]
Here the direct image \(T_{\Lambda}f(I)\) is given by
\[
    (T_{\Lambda}f)(I)(y)
    =
    \bigjoin_{x\in X}I(x)\tensor Y(y,fx).
\]

For fuzzy domains \(B\) and \(C\), the internal hom
\[
    [B,C]_{\mathrm{ff}}
\]
has as objects the fuzzy-domain morphisms
\[
    B\to C.
\]
Its hom-values are inherited from the enriched functor preorder:
\[
    [B,C]_{\mathrm{ff}}(f,g)
    =
    \bigmeet_{b\in B}C(fb,gb).
\]
Finite-flat joins in \([B,C]_{\mathrm{ff}}\) are computed pointwise in \(C\).

For \(\Lambda=2\), finite-flat weights are ordinary ideals.  On separated
preorders, fuzzy domains are precisely dcpos, and fuzzy-domain morphisms are
precisely Scott-continuous maps.  Hence
\[
    \mathbf{FDom}_{2}^{\mathrm{sep}}
    \simeq
    \mathbf{DCPO},
\]
where \(\mathbf{DCPO}\) denotes the category of dcpos and Scott-continuous maps,
small in the chosen universe.

In this crisp separated case, finite-flat bimorphisms are monotone maps
\[
    A\times B\to C
\]
between dcpos that are Scott-continuous in each variable separately.  Such maps
are Scott-continuous for the product dcpo.  Indeed, let
\[
    D\subseteq A\times B
\]
be nonempty and directed.  Then \(\pi_A D\) and \(\pi_BD\) are nonempty directed
subsets, and
\[
    \bigjoin D
    =
    \left(
        \bigjoin\pi_A D,
        \bigjoin\pi_BD
    \right).
\]
If
\[
    f:A\times B\to C
\]
is separately Scott-continuous, then
\[
\begin{aligned}
    f\left(\bigjoin D\right)
    &=
    f\left(
        \bigjoin\pi_A D,
        \bigjoin\pi_B D
    \right)
    \\
    &=
    \bigjoin_{a\in\pi_A D}
        f\left(a,\bigjoin\pi_BD\right)
    \\
    &=
    \bigjoin_{a\in\pi_A D}
        \bigjoin_{b\in\pi_BD}
            f(a,b).
\end{aligned}
\]
Since \(D\) is directed, every pair
\[
    (a,b)\in\pi_A D\times\pi_BD
\]
is bounded above by some element of \(D\).  By monotonicity,
\[
    \bigjoin_{a\in\pi_A D,\;b\in\pi_BD}f(a,b)
    =
    \bigjoin_{(a,b)\in D}f(a,b).
\]
Hence
\[
    f\left(\bigjoin D\right)
    =
    \bigjoin_{(a,b)\in D}f(a,b).
\]
Thus \(f\) is Scott-continuous on the product dcpo.

Consequently, in the separated crisp specialization, the Kelly tensor product
\(\boxtimes_{\mathrm{ff}}\) agrees with the cartesian product of dcpos, and the
internal hom agrees with the dcpo of Scott-continuous maps ordered pointwise.
Hence, under the equivalence
\[
    \mathbf{FDom}_{2}^{\mathrm{sep}}
    \simeq
    \mathbf{DCPO},
\]
the symmetric monoidal closed structure constructed above specializes to the
standard cartesian closed structure on the category of dcpos and Scott-continuous
maps.

\section{The Kleisli category and fuzzy approximable maps}
\label{sec:kleisli-fuzzy-ideal-monad}

\subsection{The enriched Kleisli category}
\label{subsec:kleisli-enriched}

\begin{definition}[Fuzzy Kleisli morphism]
\label{def:fuzzy-kleisli-morphism-enriched}
A \textbf{fuzzy Kleisli morphism}
\[
    X\rightsquigarrow Y
\]
is a \(\mathbb V\)-functor
\[
    X\to T_{\mathbb V}Y
    =
    \operatorname{Idl}_{\mathbb V}(Y).
\]
Equivalently, it is a profunctor-like object
\[
    R:X\otimes Y^{\op}\to\underline{\mathbb V}
\]
whose fibre
\[
    R(x,-):Y^{\op}\to\underline{\mathbb V}
\]
is a fuzzy ideal of \(Y\), functorially in \(x\).
\end{definition}

\begin{proposition}[Kleisli composition]
\label{prop:kleisli-composition-enriched}
If
\[
    R:X\rightsquigarrow Y
    \qquad\text{and}\qquad
    S:Y\rightsquigarrow Z
\]
are fuzzy Kleisli morphisms, their composite is
\[
    (S\circ R)(x,z)
    =
    \int^{y\in Y}R(x,y)\tensor S(y,z).
\]
The identity Kleisli morphism on \(X\) is
\[
    1_X(x,y)=X(y,x).
\]
\end{proposition}

\begin{proof}
This is Kleisli composition for the monad \(T_{\mathbb V}\), written under the
finite-flat-weight representation.  The identity is the unit
\(\eta_X(x)=X(-,x)\).
\end{proof}

\begin{proposition}[Extension to free fuzzy domains]
\label{prop:kleisli-extension-free-domain}
A fuzzy Kleisli morphism
\[
    R:X\rightsquigarrow Y
\]
extends to a morphism between free fuzzy domains
\[
    \overline R:
    \operatorname{Idl}_{\mathbb V}(X)
    \to
    \operatorname{Idl}_{\mathbb V}(Y)
\]
given by
\[
    \overline R(I)(y)
    =
    \int^{x\in X} I(x)\tensor R(x,y).
\]
\end{proposition}

\begin{proof}
This is the Kleisli extension
\[
    T_{\mathbb V}X
    \xrightarrow{T_{\mathbb V}R}
    T_{\mathbb V}T_{\mathbb V}Y
    \xrightarrow{\mu_Y}
    T_{\mathbb V}Y.
\]
The displayed formula is the resulting convolution.
\end{proof}

\subsection{The posetal specialization}
\label{subsec:kleisli-posetal}

\begin{definition}[Fuzzy approximable map]
\label{def:fuzzy-approximable-map}
For fuzzy preorders \(X,Y\), a \textbf{fuzzy approximable map}
\[
    R:X\rightsquigarrow Y
\]
is a map
\[
    R:X\times Y\to L
\]
such that each fibre
\[
    R(x,-):Y^{\op}\to\Lambda
\]
is a fuzzy ideal of \(Y\), and the assignment
\[
    x\longmapsto R(x,-)
\]
is a \(\Lambda\)-functor
\[
    X\to\operatorname{Idl}_\Lambda(Y).
\]
\end{definition}

\begin{proposition}[Posetal Kleisli formulas]
\label{prop:posetal-kleisli-formulas}
For fuzzy approximable maps
\[
    R:X\rightsquigarrow Y,
    \qquad
    S:Y\rightsquigarrow Z,
\]
composition is
\[
    (S\circ R)(x,z)
    =
    \bigjoin_{y\in Y}R(x,y)\tensor S(y,z).
\]
The identity is
\[
    1_X(x,y)=X(y,x).
\]
The extension of \(R\) to free fuzzy domains is
\[
    \overline R(I)(y)
    =
    \bigjoin_{x\in X}I(x)\tensor R(x,y).
\]
\end{proposition}

\begin{proof}
These are the posetal specializations of
Propositions~\ref{prop:kleisli-composition-enriched} and
\ref{prop:kleisli-extension-free-domain}.
\end{proof}

\begin{example}[The crisp Kleisli category]
\label{ex:crisp-kleisli}
For \(\Lambda=2\), a Kleisli morphism
\[
    X\rightsquigarrow Y
\]
is a monotone map
\[
    X\to\operatorname{Idl}(Y).
\]
Equivalently, it is an ideal-valued relation
\[
    R\subseteq X\times Y
\]
such that each fibre
\[
    R[x]=\{y\in Y\mid xRy\}
\]
is an ideal of \(Y\), and
\[
    x\leq x'
    \quad\Longrightarrow\quad
    R[x]\subseteq R[x'].
\]
Composition is
\[
    (S\circ R)[x]
    =
    \bigcup_{y\in R[x]}S[y],
\]
and the Kleisli extension is
\[
    \overline r(I)=\bigcup_{x\in I}r(x).
\]
These are the classical approximable maps between bases.
\end{example}

\begin{example}[Lawvere Kleisli morphisms]
\label{ex:lawvere-kleisli}
For the Lawvere quantale, Kleisli morphisms are flat-weight-valued metric
relations.  Composition is infimal convolution:
\[
    (S\circ R)(x,z)
    =
    \inf_{y\in Y}\bigl(R(x,y)+S(y,z)\bigr).
\]
\end{example}

\section{Fuzzy algebraic and fuzzy continuous domains}
\label{sec:fuzzy-algebraic-continuous-domains}

\subsection{Fuzzy algebraic domains}
\label{subsec:fuzzy-algebraic-domains}

\begin{definition}[Presented fuzzy algebraic domain]
\label{def:presented-fuzzy-algebraic-domain}
A \textbf{presented fuzzy algebraic domain} is a free fuzzy domain
\[
    \operatorname{Idl}_{\mathbb V}(B)
\]
for a small \(\mathbb V\)-category \(B\).  The objects of \(B\) are regarded as
basic elements, and the principal ideals
\[
    \eta_B(b)=B(-,b)
\]
are their basic compact approximants.
\end{definition}

\begin{proposition}[Algebra structure on a free fuzzy domain]
\label{prop:free-fuzzy-domain-algebra-structure}
The algebra structure on
\[
    \operatorname{Idl}_{\mathbb V}(B)
\]
is the monad multiplication
\[
    \mu_B:
    \operatorname{Idl}_{\mathbb V}(\operatorname{Idl}_{\mathbb V}(B))
    \to
    \operatorname{Idl}_{\mathbb V}(B).
\]
In the posetal case,
\[
    \mu_B(\mathcal I)(b)
    =
    \bigjoin_{I\in\operatorname{Idl}_\Lambda(B)}
    \mathcal I(I)\tensor I(b).
\]
\end{proposition}

\begin{proof}
This is the free-algebra construction for a monad, with multiplication given by
flattening.
\end{proof}

\begin{example}[Crisp algebraic domains]
\label{ex:crisp-algebraic-domains}
For \(\Lambda=2\),
\[
    \operatorname{Idl}(B)
\]
is an algebraic dcpo with compact elements the principal ideals
\[
    \downarrow b.
\]
Every ideal is the directed supremum of its principal subideals:
\[
    I=\bigjoin_{b\in I}\downarrow b.
\]
Thus free algebras for the crisp ideal monad are algebraic dcpos equipped with a
chosen basis of compact generators.
\end{example}

\subsection{Fuzzy continuous domains}
\label{subsec:fuzzy-continuous-domains}

\begin{definition}[Fuzzy continuous domain]
\label{def:fuzzy-continuous-domain-enriched}
Let \(X\) be a fuzzy domain with algebra structure
\[
    \alpha:\operatorname{Idl}_{\mathbb V}(X)\to X.
\]
We say that \(X\) is a \textbf{fuzzy continuous domain} if \(\alpha\) admits an
enriched left adjoint
\[
    \beta:X\to\operatorname{Idl}_{\mathbb V}(X),
    \qquad
    \beta\dashv\alpha.
\]
Thus
\[
    \operatorname{Idl}_{\mathbb V}(X)(\beta x,I)
    \cong
    X(x,\alpha I).
\]
The object \(\beta(x)\) is the fuzzy ideal of approximants to \(x\).
\end{definition}

\begin{definition}[Fuzzy way-below degree]
\label{def:fuzzy-way-below-degree}
In the posetal specialization, the associated \textbf{fuzzy way-below degree} is
\[
    \operatorname{wb}_X(y,x)
    =
    \beta(x)(y).
\]
\end{definition}

\begin{proposition}[Adjunction and approximation]
\label{prop:fuzzy-continuity-adjunction}
Let \(X\) be a separated fuzzy domain over \(\Lambda\), with algebra map
\[
    \alpha:\operatorname{Idl}_\Lambda(X)\to X.
\]
If
\[
    \beta:X\to\operatorname{Idl}_\Lambda(X)
\]
satisfies
\[
    \beta\dashv\alpha,
\]
then
\[
    \alpha\beta=1_X
\]
in the induced crisp order.  Hence every element is recovered as the fuzzy join of
its approximants.
\end{proposition}

\begin{proof}
The unit of the adjunction gives
\[
    1_X\leq \alpha\beta.
\]
The counit gives
\[
    \beta\alpha\leq 1_{\operatorname{Idl}_\Lambda(X)}.
\]
Apply the counit inequality to the principal ideal \(\eta_X(x)\).  Since
\[
    \alpha\eta_X(x)=x,
\]
we obtain
\[
    \beta(x)\leq\eta_X(x).
\]
Applying \(\alpha\) gives
\[
    \alpha\beta(x)\leq\alpha\eta_X(x)=x.
\]
Together with \(x\leq\alpha\beta(x)\), this gives equality in the induced crisp
order.  Since \(X\) is separated, this is equality of objects.
\end{proof}

\begin{example}[Crisp continuous domains]
\label{ex:crisp-continuous-domains}
Let \(D\) be a dcpo with ideal-supremum map
\[
    \bigjoin:\operatorname{Idl}(D)\to D.
\]
Then \(D\) is continuous precisely when this map has a left adjoint
\[
    \beta:D\to\operatorname{Idl}(D).
\]
The left adjoint is
\[
    \beta(x)=\{y\in D\mid y\ll x\},
\]
and the adjunction says
\[
    \beta(x)\subseteq I
    \quad\Longleftrightarrow\quad
    x\leq \bigjoin I.
\]
\end{example}

\begin{example}[Split Kleisli idempotents]
\label{ex:split-kleisli-idempotents-continuous}
In the crisp case, the map
\[
    \beta:D\to\operatorname{Idl}(D),
    \qquad
    \beta(x)=\{y\mid y\ll x\},
\]
is a Kleisli endomorphism
\[
    D\rightsquigarrow D.
\]
Interpolation of the way-below relation gives
\[
    \beta\circ_{\operatorname{Kl}}\beta=\beta.
\]
Thus continuous dcpos appear as splittings of Kleisli idempotents, equivalently as
retracts of algebraic dcpos.
\end{example}

\begin{example}[Lawvere-valued approximation]
\label{ex:lawvere-fuzzy-continuous}
For the Lawvere quantale, a fuzzy continuity structure assigns to each point \(x\)
a finite-flat weight
\[
    \beta(x):X^{\op}\to[0,\infty].
\]
The value \(\beta(x)(y)\) is the cost by which \(y\) approximates \(x\).  The
adjunction
\[
    \beta\dashv\alpha
\]
says that approximation into \(x\) is controlled exactly by comparison with
finite-flat-weight colimits.
\end{example}

\section{Fuzzy Scott opens and fuzzy Scott locales}
\label{sec:fuzzy-scott-opens}

\subsection{The enriched construction}
\label{subsec:fuzzy-scott-opens-enriched}

\begin{definition}[Fuzzy Scott open]
\label{def:fuzzy-scott-open-enriched}
Let \(X\) be a weak fuzzy domain with structure map
\[
    \alpha:T_{\mathbb V}X\to X.
\]
A \textbf{fuzzy Scott open} on \(X\) is an upper fuzzy predicate
\[
    U:X\to\underline{\mathbb V}
\]
such that, for every fuzzy ideal \(I\),
\[
    U(\alpha(I))
    \cong
    \int^{x\in X} I(x)\tensor U(x).
\]
\end{definition}

\begin{definition}[Enriched fuzzy Scott frame]
\label{def:enriched-fuzzy-scott-frame}
When the required equalizer exists in categorical fuzzy frames, define
\[
    \operatorname{Scott}_{\mathbb V}(X)
\]
to be the equalizer
\[
\begin{tikzcd}[column sep=large]
\operatorname{Scott}_{\mathbb V}(X)
    \arrow[r, hook]
&
P_{\mathbb V}(X^{\op})
    \arrow[r, shift left=.8ex, "{U\mapsto U\alpha}"]
    \arrow[r, shift right=.8ex, "{U\mapsto (I\mapsto \int^x I(x)\tensor U(x))}"']
&
P_{\mathbb V}(\operatorname{Idl}_{\mathbb V}(X)^{\op}).
\end{tikzcd}
\]
The second map is a fuzzy-frame morphism because, for each finite-flat weight \(I\),
the functional
\[
    U\longmapsto I*U
\]
is a character, and functoriality in \(I\) gives an upper predicate on
\(\operatorname{Idl}_{\mathbb V}(X)\).
\end{definition}

\begin{definition}[Fuzzy Scott locale]
\label{def:fuzzy-scott-locale-enriched}
The \textbf{fuzzy Scott locale} of \(X\) is the fuzzy locale
\[
    \operatorname{ScottLoc}_{\mathbb V}(X)
\]
whose frame of opens is
\[
    \operatorname{Opens}_{\mathbb V}(\operatorname{ScottLoc}_{\mathbb V}(X))
    =
    \operatorname{Scott}_{\mathbb V}(X).
\]
The inclusion
\[
    \operatorname{Scott}_{\mathbb V}(X)
    \hookrightarrow
    P_{\mathbb V}(X^{\op})
\]
defines a fuzzy-locale map
\[
    \mathcal P_{\mathbb V}^{\uparrow}X
    \longrightarrow
    \operatorname{ScottLoc}_{\mathbb V}(X).
\]
\end{definition}

\subsection{The posetal specialization}
\label{subsec:fuzzy-scott-opens-posetal}

\begin{definition}[Fuzzy Scott open over \(\Lambda\)]
\label{def:fuzzy-scott-open-posetal}
Let \(X\) be a weak fuzzy domain over \(\Lambda\), with structure map
\[
    \alpha:\operatorname{Idl}_\Lambda(X)\to X.
\]
A \textbf{fuzzy Scott open} is an upper fuzzy predicate
\[
    U:X\to L
\]
satisfying
\[
    U(\alpha(I))
    =
    \bigjoin_{x\in X} I(x)\tensor U(x)
\]
for every fuzzy ideal \(I\).
\end{definition}

\begin{definition}[The fuzzy Scott frame]
\label{def:fuzzy-scott-frame-posetal}
The \textbf{fuzzy Scott frame} of \(X\) is the full fuzzy subpreorder
\[
    \operatorname{Scott}_\Lambda(X)
    \hookrightarrow
    P_\Lambda(X^{\op})
\]
spanned by the fuzzy Scott opens.  Equivalently, it is the equalizer
\[
\begin{tikzcd}[column sep=large]
\operatorname{Scott}_\Lambda(X)
    \arrow[r, hook]
&
P_\Lambda(X^{\op})
    \arrow[r, shift left=.8ex, "{U\mapsto U\alpha}"]
    \arrow[r, shift right=.8ex, "{\Theta}"']
&
P_\Lambda(\operatorname{Idl}_\Lambda(X)^{\op}),
\end{tikzcd}
\]
where
\[
    \Theta(U)(I)
    =
    \bigjoin_{x\in X} I(x)\tensor U(x).
\]
\end{definition}

\begin{proposition}[Fuzzy Scott opens form a fuzzy frame]
\label{prop:fuzzy-scott-opens-frame}
Let \(X\) be a weak fuzzy domain over \(\Lambda\).  Then
\[
    \operatorname{Scott}_\Lambda(X)
\]
is a fuzzy frame, and the inclusion
\[
    \operatorname{Scott}_\Lambda(X)\hookrightarrow P_\Lambda(X^{\op})
\]
is a fuzzy-frame homomorphism.
\end{proposition}

\begin{proof}
First, \(\Theta(U)\) is an upper fuzzy predicate on
\(\operatorname{Idl}_\Lambda(X)\).  If
\[
    r\leq \operatorname{Idl}_\Lambda(X)(I,J),
\]
then \(r\tensor I(x)\leq J(x)\) for every \(x\), and hence
\[
    r\tensor \Theta(U)(I)
    =
    r\tensor \bigjoin_x I(x)\tensor U(x)
    \leq
    \bigjoin_x J(x)\tensor U(x)
    =
    \Theta(U)(J).
\]
Thus
\[
    \Theta(U)\in
    P_\Lambda(\operatorname{Idl}_\Lambda(X)^{\op}).
\]

For each fuzzy ideal \(I\), the map
\[
    U\longmapsto \Theta(U)(I)
\]
is the character \(\chi_I\), so it preserves arbitrary joins, finite meets, and
scalars.  Since fuzzy-frame structure in
\[
    P_\Lambda(\operatorname{Idl}_\Lambda(X)^{\op})
\]
is computed pointwise in \(I\), the map
\[
    \Theta:
    P_\Lambda(X^{\op})
    \to
    P_\Lambda(\operatorname{Idl}_\Lambda(X)^{\op})
\]
preserves arbitrary joins, finite meets, and scalars.  The map
\(U\mapsto U\alpha\) is precomposition by \(\alpha\), so it preserves the same
structure.  Therefore their equalizer is closed under arbitrary joins, finite
meets, and scalars, and the inclusion is a fuzzy-frame homomorphism.
\end{proof}

\begin{definition}[Fuzzy Scott locale over \(\Lambda\)]
\label{def:fuzzy-scott-locale-posetal}
For a weak fuzzy domain \(X\), the \textbf{fuzzy Scott locale} of \(X\) is the
fuzzy locale
\[
    \operatorname{ScottLoc}_\Lambda(X)
\]
whose frame of opens is
\[
    \Opens_\Lambda(\operatorname{ScottLoc}_\Lambda(X))
    =
    \operatorname{Scott}_\Lambda(X).
\]
The inclusion
\[
    \operatorname{Scott}_\Lambda(X)\hookrightarrow P_\Lambda(X^{\op})
\]
gives a fuzzy-locale map
\[
    \mathcal P_\Lambda^\uparrow X
    \longrightarrow
    \operatorname{ScottLoc}_\Lambda(X).
\]
\end{definition}

\begin{theorem}[Scott opens of a free fuzzy domain]
\label{thm:scott-opens-free-fuzzy-domain}
For every fuzzy preorder \(B\), there is a canonical isomorphism of fuzzy frames
\[
    P_\Lambda(B^{\op})
    \cong
    \operatorname{Scott}_\Lambda(\operatorname{Idl}_\Lambda(B)).
\]
An upper fuzzy predicate
\[
    U:B\to L
\]
corresponds to the fuzzy Scott open
\[
    \widehat U:\operatorname{Idl}_\Lambda(B)\to L
\]
defined by
\[
    \widehat U(I)
    =
    \bigjoin_{b\in B}I(b)\tensor U(b).
\]
\end{theorem}

\begin{proof}
Let \(U\in P_\Lambda(B^{\op})\).  Define
\[
    \widehat U(I)
    =
    \bigjoin_b I(b)\tensor U(b).
\]
The calculation in Proposition~\ref{prop:fuzzy-scott-opens-frame} shows that
\(\widehat U\) is an upper fuzzy predicate on
\(\operatorname{Idl}_\Lambda(B)\).

For
\[
    \mathcal I\in
    \operatorname{Idl}_\Lambda(\operatorname{Idl}_\Lambda(B)),
\]
we compute
\[
\begin{aligned}
    \widehat U(\mu_B(\mathcal I))
    &=
    \bigjoin_b \mu_B(\mathcal I)(b)\tensor U(b)\\
    &=
    \bigjoin_b
    \left(
        \bigjoin_I \mathcal I(I)\tensor I(b)
    \right)\tensor U(b)\\
    &=
    \bigjoin_I \mathcal I(I)\tensor
    \left(
        \bigjoin_b I(b)\tensor U(b)
    \right)\\
    &=
    \bigjoin_I \mathcal I(I)\tensor \widehat U(I).
\end{aligned}
\]
Thus \(\widehat U\) is Scott open.

Conversely, let
\[
    W\in\operatorname{Scott}_\Lambda(\operatorname{Idl}_\Lambda(B)).
\]
Restrict along the unit and define
\[
    U_W(b)=W(\eta_B b).
\]
Then \(U_W\) is an upper fuzzy predicate on \(B\).  For \(U\), Yoneda gives
\[
    \widehat U(\eta_B b)
    =
    \bigjoin_c B(c,b)\tensor U(c)
    =
    U(b),
\]
so \(U_{\widehat U}=U\).

It remains to show that \(\widehat{U_W}=W\).  For
\(I\in\operatorname{Idl}_\Lambda(B)\), the upper-predicate law for \(W\) gives
\[
    \operatorname{Idl}_\Lambda(B)(\eta_B b,I)
    \tensor W(\eta_B b)
    \leq
    W(I).
\]
By Yoneda,
\[
    \operatorname{Idl}_\Lambda(B)(\eta_B b,I)=I(b).
\]
Hence
\[
    \widehat{U_W}(I)
    =
    \bigjoin_b I(b)\tensor W(\eta_B b)
    \leq
    W(I).
\]

For the converse, use
\[
    \mu_B T_\Lambda(\eta_B)=1_{\operatorname{Idl}_\Lambda(B)}.
\]
Since \(W\) is Scott open,
\[
\begin{aligned}
    W(I)
    &=
    W(\mu_B T_\Lambda(\eta_B)(I))\\
    &=
    \bigjoin_J T_\Lambda(\eta_B)(I)(J)\tensor W(J).
\end{aligned}
\]
Moreover,
\[
    T_\Lambda(\eta_B)(I)(J)
    =
    \bigjoin_b I(b)\tensor
    \operatorname{Idl}_\Lambda(B)(J,\eta_B b).
\]
Using the upper-predicate law for \(W\),
\[
    \operatorname{Idl}_\Lambda(B)(J,\eta_B b)\tensor W(J)
    \leq
    W(\eta_B b).
\]
Therefore
\[
\begin{aligned}
    W(I)
    &\leq
    \bigjoin_J \bigjoin_b
    I(b)\tensor
    \operatorname{Idl}_\Lambda(B)(J,\eta_B b)\tensor W(J)\\
    &\leq
    \bigjoin_b I(b)\tensor W(\eta_B b)\\
    &=
    \widehat{U_W}(I).
\end{aligned}
\]
Thus \(W=\widehat{U_W}\).  The two constructions are inverse and preserve the
fuzzy-frame structure because all operations are computed pointwise and each
evaluation functional is a character.
\end{proof}

\begin{example}[The crisp Scott condition]
\label{ex:crisp-scott-opens}
For \(\Lambda=2\), if \(X\) is a dcpo with algebra
\[
    \alpha:\operatorname{Idl}(X)\to X,
    \qquad
    I\mapsto\bigjoin I,
\]
then a fuzzy Scott open is an upper set \(U\subseteq X\) satisfying
\[
    \bigjoin I\in U
    \quad\Longleftrightarrow\quad
    I\cap U\neq\varnothing
\]
for every ideal \(I\).  This is the ordinary Scott-open condition expressed using
ideals.
\end{example}

\begin{example}[Upper sets as Scott opens of ideal completions]
\label{ex:upper-sets-as-scott-opens-ideal-completion}
For a preorder \(B\),
\[
    \operatorname{Scott}(\operatorname{Idl}(B))
    \cong
    \operatorname{Up}(B).
\]
The correspondence sends an upper set \(U\subseteq B\) to
\[
    \widehat U
    =
    \{I\in\operatorname{Idl}(B)\mid I\cap U\neq\varnothing\}.
\]
Thus the upper Alexandrov frame of \(B\) is the Scott frame of the free dcpo
generated by \(B\).
\end{example}

\begin{example}[Standard fuzzy Scott formulas]
\label{ex:standard-fuzzy-scott-opens}
For G\"odel truth values, the fuzzy Scott condition is
\[
    U(\alpha(I))
    =
    \bigjoin_x \min(I(x),U(x)).
\]
For {\L}ukasiewicz truth values, it is
\[
    U(\alpha(I))
    =
    \bigjoin_x \max(0,I(x)+U(x)-1).
\]
For product truth values, it is
\[
    U(\alpha(I))
    =
    \bigjoin_x I(x)U(x).
\]
For the Lawvere quantale, it becomes
\[
    U(\alpha(I))
    =
    \inf_x\bigl(I(x)+U(x)\bigr).
\]
\end{example}

\section{The Alexandrov character comonad}
\label{sec:alexandrov-character-comonad}

\subsection{The enriched comonad}
\label{subsec:alexandrov-character-comonad-enriched}

\begin{definition}[Alexandrov character comonad]
\label{def:alexandrov-character-comonad-enriched}
The \textbf{Alexandrov character comonad} is the comonad
\[
    G_{\mathbb V}
    =
    \mathcal P_{\mathbb V}^{\uparrow}\operatorname{Char}_{\mathbb V}
\]
on fuzzy locales induced by
\[
    \mathcal P_{\mathbb V}^{\uparrow}
    \dashv
    \operatorname{Char}_{\mathbb V}.
\]
For a fuzzy locale \(\mathcal X\),
\[
    G_{\mathbb V}\mathcal X
    =
    \mathcal P_{\mathbb V}^{\uparrow}
    \bigl(
        \operatorname{Char}_{\mathbb V}(\mathcal X)
    \bigr),
\]
so
\[
    \operatorname{Opens}_{\mathbb V}(G_{\mathbb V}\mathcal X)
    =
    P_{\mathbb V}
    \bigl(
        \operatorname{Char}_{\mathbb V}(\mathcal X)^{\op}
    \bigr).
\]
\end{definition}

\begin{proposition}[Counit of the character comonad]
\label{prop:alexandrov-character-comonad-counit}
The counit
\[
    G_{\mathbb V}\mathcal X\to\mathcal X
\]
is frame-dual to the evaluation fuzzy-frame morphism
\[
    \operatorname{Opens}_{\mathbb V}(\mathcal X)
    \longrightarrow
    P_{\mathbb V}
    \bigl(
        \operatorname{Char}_{\mathbb V}(\mathcal X)^{\op}
    \bigr),
\]
given by
\[
    a\longmapsto \widehat a,
    \qquad
    \widehat a(\chi)=\chi(a).
\]
\end{proposition}

\begin{proof}
This is the counit of the adjunction, evaluated at \(\mathcal X\).  On frames it is
the transpose of the identity functor on
\(\operatorname{Char}_{\mathbb V}(\mathcal X)\), namely evaluation.
\end{proof}

\begin{definition}[Coalgebra for the Alexandrov character comonad]
\label{def:alexandrov-character-coalgebra}
A \textbf{\(G_{\mathbb V}\)-coalgebra} is a fuzzy locale \(\mathcal X\) equipped
with a fuzzy-locale map
\[
    \mathcal X\to G_{\mathbb V}\mathcal X
\]
satisfying the usual comonad coalgebra laws.  Frame-dually, this says that
\[
    \operatorname{Opens}_{\mathbb V}(\mathcal X)
\]
is a comonad-coherent retract of the upper fuzzy predicate frame on its character
object.
\end{definition}

\subsection{The posetal specialization}
\label{subsec:alexandrov-character-comonad-posetal}

\begin{proposition}[The posetal character comonad]
\label{prop:posetal-character-comonad}
For a fuzzy locale \(\mathcal X\),
\[
    \Opens_\Lambda(G_\Lambda\mathcal X)
    =
    P_\Lambda
    \bigl(
        \operatorname{Char}_\Lambda(\Opens_\Lambda(\mathcal X))^{\op}
    \bigr).
\]
The counit is frame-dual to the evaluation map
\[
    a\longmapsto \widehat a,
    \qquad
    \widehat a(\chi)=\chi(a).
\]
A \(G_\Lambda\)-coalgebra is frame-dually a fuzzy-frame homomorphism
\[
    P_\Lambda
    \bigl(
        \operatorname{Char}_\Lambda(\Opens_\Lambda(\mathcal X))^{\op}
    \bigr)
    \to
    \Opens_\Lambda(\mathcal X)
\]
compatible with the counit and comultiplication.
\end{proposition}

\begin{proof}
This is the posetal specialization of the enriched comonad and its counit.
\end{proof}

\begin{example}[The crisp case]
\label{ex:crisp-character-comonad}
For \(\Lambda=2\), the character object of a locale is its preorder of points,
ordered by specialization according to the convention
\[
    p\leq q
    \quad\Longleftrightarrow\quad
    \forall a,\;p(a)\leq q(a).
\]
Thus
\[
    G_2\mathcal X
\]
is the upper Alexandrov locale of the specialization preorder of points of
\(\mathcal X\).  The counit is frame-dual to
\[
    \Opens(\mathcal X)
    \to
    \operatorname{Up}(\operatorname{Pt}(\mathcal X)),
    \qquad
    a\longmapsto \{p\mid p(a)=\top\}.
\]
\end{example}

\begin{example}[Presheaf fuzzy locales]
\label{ex:presheaf-fuzzy-locale-character-comonad}
If
\[
    \mathcal X=\mathcal P_\Lambda^{\uparrow}X,
\]
then
\[
    \operatorname{Char}_\Lambda(\Opens_\Lambda(\mathcal X))
    =
    \operatorname{Char}_\Lambda(P_\Lambda(X^{\op}))
    \cong
    \operatorname{Idl}_\Lambda(X).
\]
Therefore
\[
    G_\Lambda(\mathcal P_\Lambda^{\uparrow}X)
\]
has opens given by upper fuzzy predicates on
\(\operatorname{Idl}_\Lambda(X)\).
\end{example}

\section{The co-Kleisli category}
\label{sec:cokleisli-character-comonad}

\subsection{The enriched co-Kleisli category}
\label{subsec:cokleisli-enriched}

\begin{theorem}[Co-Kleisli morphisms]
\label{thm:cokleisli-morphisms-enriched}
For fuzzy locales \(\mathcal X,\mathcal Y\),
\[
    \operatorname{CoKl}(G_{\mathbb V})(\mathcal X,\mathcal Y)
    =
    \mathbf{FLoc}_{\mathbb V}(G_{\mathbb V}\mathcal X,\mathcal Y)
\]
is naturally equivalent to
\[
    \mathbb V\text{-}\mathbf{Cat}
    \bigl(
        \operatorname{Char}_{\mathbb V}(\mathcal X),
        \operatorname{Char}_{\mathbb V}(\mathcal Y)
    \bigr).
\]
Thus a co-Kleisli morphism
\[
    \mathcal X\to\mathcal Y
\]
is equivalently a \(\mathbb V\)-functor between the character objects.
\end{theorem}

\begin{proof}
By definition,
\[
    G_{\mathbb V}\mathcal X
    =
    \mathcal P_{\mathbb V}^{\uparrow}
    \operatorname{Char}_{\mathbb V}(\mathcal X).
\]
The upper fuzzy character adjunction gives
\[
\begin{aligned}
    \mathbf{FLoc}_{\mathbb V}(G_{\mathbb V}\mathcal X,\mathcal Y)
    &=
    \mathbf{FLoc}_{\mathbb V}
    \bigl(
        \mathcal P_{\mathbb V}^{\uparrow}
        \operatorname{Char}_{\mathbb V}(\mathcal X),
        \mathcal Y
    \bigr)\\
    &\simeq
    \mathbb V\text{-}\mathbf{Cat}
    \bigl(
        \operatorname{Char}_{\mathbb V}(\mathcal X),
        \operatorname{Char}_{\mathbb V}(\mathcal Y)
    \bigr).
\end{aligned}
\]
\end{proof}

\subsection{The posetal specialization}
\label{subsec:cokleisli-posetal}

\begin{corollary}[Posetal co-Kleisli morphisms]
\label{cor:posetal-cokleisli-morphisms}
In the posetal specialization, a co-Kleisli morphism
\[
    \mathcal X\to\mathcal Y
\]
is a \(\Lambda\)-functor
\[
    \operatorname{Char}_\Lambda(\Opens_\Lambda(\mathcal X))
    \to
    \operatorname{Char}_\Lambda(\Opens_\Lambda(\mathcal Y)).
\]
Such a map need not arise from a fuzzy-locale map
\[
    \mathcal X\to\mathcal Y.
\]
\end{corollary}

\begin{proof}
This is the posetal specialization of
Theorem~\ref{thm:cokleisli-morphisms-enriched}.
\end{proof}

\begin{example}[The crisp case]
\label{ex:crisp-cokleisli}
For \(\Lambda=2\), co-Kleisli morphisms are monotone maps between the specialization
preorders of points:
\[
    \operatorname{Pt}(\mathcal X)
    \to
    \operatorname{Pt}(\mathcal Y).
\]
These need not be continuous maps of locales or spaces.
\end{example}

\begin{example}[Presheaf fuzzy locales]
\label{ex:presheaf-fuzzy-locale-cokleisli}
Between upper presheaf fuzzy locales, co-Kleisli morphisms are fuzzy functors
between fuzzy ideal completions:
\[
    \operatorname{Idl}_\Lambda(X)\to\operatorname{Idl}_\Lambda(Y).
\]
\end{example}

\begin{example}[Lawvere co-Kleisli maps]
\label{ex:lawvere-cokleisli}
For the Lawvere quantale, co-Kleisli morphisms are nonexpansive maps between metric
spaces of finite-flat weights.
\end{example}

\section{Summary}
\label{sec:summary-fuzzy-ideal-monad-domains-scott}

The upper fuzzy predicate frame of \(X\) is
\[
    P_{\mathbb V}(X^{\op})
    =
    [X,\underline{\mathbb V}],
\]
and in the posetal specialization it is
\[
    P_\Lambda(X^{\op})
    =
    [X,\Lambda].
\]
The character adjunction
\[
    \mathcal P_{\mathbb V}^{\uparrow}
    \dashv
    \operatorname{Char}_{\mathbb V}
\]
induces the fuzzy ideal monad
\[
    T_\Lambda(X)
    =
    \operatorname{Idl}_\Lambda(X),
\]
whose unit is the principal-ideal embedding and whose multiplication is fuzzy
flattening.

Weak fuzzy domains carry chosen candidate fuzzy ideal-joins, while fuzzy domains
are strict algebras for this monad.  Free fuzzy domains
\[
    \operatorname{Idl}_\Lambda(B)
\]
are the presented fuzzy algebraic domains.  Fuzzy continuous domains are those for
which
\[
    \alpha:\operatorname{Idl}_\Lambda(X)\to X
\]
has an enriched left adjoint
\[
    \beta:X\to\operatorname{Idl}_\Lambda(X),
\]
with fuzzy way-below degree
\[
    \operatorname{wb}_X(y,x)=\beta(x)(y).
\]

Kleisli morphisms for the fuzzy ideal monad are fuzzy ideal-valued maps
\[
    X\rightsquigarrow Y,
\]
or fuzzy approximable maps.  Fuzzy Scott opens of a weak fuzzy domain
\[
    \alpha:\operatorname{Idl}_\Lambda(X)\to X
\]
are exactly the upper fuzzy predicates \(U\) satisfying
\[
    U(\alpha(I))
    =
    \bigjoin_x I(x)\tensor U(x).
\]
For \(\Lambda=2\), these constructions recover ordinary ideals, ideal completion,
dcpos, algebraic dcpos, continuous dcpos, approximable maps, and ordinary Scott
opens.